\documentclass[11pt]{article}
% =======================================================================
%   Shared style for CS 300 notes.
%   Usage: place `% =======================================================================
%   Shared style for CS 300 notes.
%   Usage: place `% =======================================================================
%   Shared style for CS 300 notes.
%   Usage: place `\input{../../notes-style.tex}` right after \documentclass
%   in each chapter's lectures.tex and notes.tex.
% =======================================================================

% ---- Layout -----------------------------------------------------------
\usepackage[margin=1in]{geometry}
\usepackage{parskip}        % space between paragraphs, no indent
\usepackage{microtype}      % subtle kerning/spacing improvements

% ---- Colors (muted palette, easy on light-sensitive eyes) -------------
\usepackage{xcolor}
\definecolor{pagebg}       {HTML}{FAF6F0}  % warm cream background
\definecolor{bodyfg}       {HTML}{3D3A36}  % warm dark gray body text
\definecolor{heading}      {HTML}{3D6A6B}  % muted teal
\definecolor{subheading}   {HTML}{5B7A9C}  % dusty blue
\definecolor{accent}       {HTML}{8A6B4A}  % warm tan
\definecolor{linkcol}      {HTML}{5B7A9C}  % same as subheading
\definecolor{mutedtext}    {HTML}{6B6560}  % secondary / caption text
\definecolor{rulecol}      {HTML}{C8BFB0}  % separator / rule

% Callout box palette
\definecolor{defnFrame}    {HTML}{9DB89B}  \definecolor{defnBack}    {HTML}{EEF2E8}
\definecolor{ideaFrame}    {HTML}{9FB4CD}  \definecolor{ideaBack}    {HTML}{E8EEF5}
\definecolor{warnFrame}    {HTML}{CFAE87}  \definecolor{warnBack}    {HTML}{F5EDDF}
\definecolor{noteFrame}    {HTML}{BFA6AE}  \definecolor{noteBack}    {HTML}{F2E8EE}
\definecolor{exFrame}      {HTML}{9FBDBD}  \definecolor{exBack}      {HTML}{E8F0F0}
\definecolor{codeBack}     {HTML}{F3EDE2}
\definecolor{codeBorder}   {HTML}{D4CBBA}
\definecolor{codeKeyword}  {HTML}{6B4F7A}
\definecolor{codeString}   {HTML}{8A6B4A}
\definecolor{codeComment}  {HTML}{8A857F}

% ---- Page background + base text color --------------------------------
\usepackage{pagecolor}
\pagecolor{pagebg}
\color{bodyfg}

% ---- Fonts ------------------------------------------------------------
\usepackage[T1]{fontenc}
\IfFileExists{libertine.sty}{%
    \usepackage{libertine}      % warm, readable serif
    \usepackage[scaled=0.95]{inconsolata}  % softer monospace
}{}

% ---- Hyperref ---------------------------------------------------------
\usepackage{hyperref}
\hypersetup{
    colorlinks=true,
    linkcolor=linkcol,
    urlcolor=linkcol,
    citecolor=linkcol,
    pdfborder={0 0 0},
}

% ---- Section styling --------------------------------------------------
\usepackage[explicit]{titlesec}
\titleformat{\section}
    {\color{heading}\Large\bfseries}
    {\color{heading}\thesection\hspace{0.6em}}{0pt}{#1}
\titleformat{\subsection}
    {\color{subheading}\large\bfseries}
    {\color{subheading}\thesubsection\hspace{0.5em}}{0pt}{#1}
\titleformat{\subsubsection}
    {\color{accent}\normalsize\bfseries}
    {\color{accent}\thesubsubsection\hspace{0.5em}}{0pt}{#1}
\titleformat{\paragraph}[runin]
    {\color{accent}\bfseries}{}{0pt}{#1.}

\titlespacing*{\section}{0pt}{1.6em}{0.6em}
\titlespacing*{\subsection}{0pt}{1.2em}{0.4em}

% ---- Enumerate / itemize spacing --------------------------------------
\usepackage{enumitem}
\setlist[itemize]{leftmargin=1.4em, itemsep=2pt, topsep=4pt}
\setlist[enumerate]{leftmargin=1.6em, itemsep=2pt, topsep=4pt}

% ---- Code listings ----------------------------------------------------
\usepackage{listings}
\lstdefinestyle{notescode}{
    backgroundcolor=\color{codeBack},
    basicstyle=\ttfamily\small\color{bodyfg},
    keywordstyle=\color{codeKeyword}\bfseries,
    stringstyle=\color{codeString},
    commentstyle=\color{codeComment}\itshape,
    frame=single,
    framesep=6pt,
    rulecolor=\color{codeBorder},
    showstringspaces=false,
    breaklines=true,
    xleftmargin=10pt,
    xrightmargin=10pt,
    aboveskip=8pt,
    belowskip=8pt,
}
\lstset{style=notescode}

% ---- Callout boxes (tcolorbox) ----------------------------------------
\usepackage[most]{tcolorbox}

\newtcolorbox{defnbox}[1][]{
    enhanced, breakable,
    colback=defnBack, colframe=defnFrame, coltitle=bodyfg,
    title={\bfseries Definition \textemdash\ #1},
    fonttitle=\bfseries,
    boxrule=0.5pt, arc=2pt, left=8pt, right=8pt, top=6pt, bottom=6pt,
    attach boxed title to top left={xshift=10pt, yshift=-6pt},
    boxed title style={colback=defnFrame, colframe=defnFrame, arc=2pt},
    before skip=8pt, after skip=8pt,
}
\newtcolorbox{ideabox}[1][]{
    enhanced, breakable,
    colback=ideaBack, colframe=ideaFrame, coltitle=bodyfg,
    title={\bfseries Key idea #1},
    boxrule=0.5pt, arc=2pt, left=8pt, right=8pt, top=6pt, bottom=6pt,
    attach boxed title to top left={xshift=10pt, yshift=-6pt},
    boxed title style={colback=ideaFrame, colframe=ideaFrame, arc=2pt},
    before skip=8pt, after skip=8pt,
}
\newtcolorbox{warnbox}[1][]{
    enhanced, breakable,
    colback=warnBack, colframe=warnFrame, coltitle=bodyfg,
    title={\bfseries Gotcha #1},
    boxrule=0.5pt, arc=2pt, left=8pt, right=8pt, top=6pt, bottom=6pt,
    attach boxed title to top left={xshift=10pt, yshift=-6pt},
    boxed title style={colback=warnFrame, colframe=warnFrame, arc=2pt},
    before skip=8pt, after skip=8pt,
}
\newtcolorbox{notebox}[1][]{
    enhanced, breakable,
    colback=noteBack, colframe=noteFrame, coltitle=bodyfg,
    title={\bfseries Aside #1},
    boxrule=0.5pt, arc=2pt, left=8pt, right=8pt, top=6pt, bottom=6pt,
    attach boxed title to top left={xshift=10pt, yshift=-6pt},
    boxed title style={colback=noteFrame, colframe=noteFrame, arc=2pt},
    before skip=8pt, after skip=8pt,
}
\newtcolorbox{examplebox}[1][]{
    enhanced, breakable,
    colback=exBack, colframe=exFrame, coltitle=bodyfg,
    title={\bfseries Example #1},
    boxrule=0.5pt, arc=2pt, left=8pt, right=8pt, top=6pt, bottom=6pt,
    attach boxed title to top left={xshift=10pt, yshift=-6pt},
    boxed title style={colback=exFrame, colframe=exFrame, arc=2pt},
    before skip=8pt, after skip=8pt,
}

% ---- TikZ graph styles (added 2026-04-25 for ch_10 augmentation) ------
% tcolorbox[most] already loads tikz; this block declares the libraries
% + shared `vertex` / `edge` styles used by ch_10's general directed-graph
% diagrams. ch_9's tree-style TikZ uses inline overrides and is
% unaffected. Simple circle nodes + arrow edges only -- no production
% styling.
\usetikzlibrary{arrows.meta, positioning, calc}
\tikzset{
    vertex/.style={circle, draw, minimum size=7mm, inner sep=0pt,
                   font=\small, thick},
    vsource/.style={vertex, fill=ideaBack},
    vdone/.style={vertex, fill=defnBack},
    edge/.style={->, >=Stealth, thick},
    uedge/.style={-, thick},
    treeedge/.style={->, >=Stealth, thick},
    backedge/.style={->, >=Stealth, thick, dashed, red!70!black},
    fwdedge/.style={->, >=Stealth, thick, blue!60!black},
    crossedge/.style={->, >=Stealth, thick, green!50!black, dotted},
    mstedge/.style={-, very thick, accent},
}

% ---- Title block ------------------------------------------------------
\usepackage{titling}
\pretitle{\begin{center}\color{heading}\LARGE\bfseries}
\posttitle{\end{center}\vskip -0.4em}
\preauthor{\begin{center}\color{mutedtext}\normalsize}
\postauthor{\end{center}}
\predate{\begin{center}\color{mutedtext}\small}
\postdate{\end{center}\vspace{0.4em}{\color{rulecol}\hrule height 0.4pt}\vspace{1.2em}}
` right after \documentclass
%   in each chapter's lectures.tex and notes.tex.
% =======================================================================

% ---- Layout -----------------------------------------------------------
\usepackage[margin=1in]{geometry}
\usepackage{parskip}        % space between paragraphs, no indent
\usepackage{microtype}      % subtle kerning/spacing improvements

% ---- Colors (muted palette, easy on light-sensitive eyes) -------------
\usepackage{xcolor}
\definecolor{pagebg}       {HTML}{FAF6F0}  % warm cream background
\definecolor{bodyfg}       {HTML}{3D3A36}  % warm dark gray body text
\definecolor{heading}      {HTML}{3D6A6B}  % muted teal
\definecolor{subheading}   {HTML}{5B7A9C}  % dusty blue
\definecolor{accent}       {HTML}{8A6B4A}  % warm tan
\definecolor{linkcol}      {HTML}{5B7A9C}  % same as subheading
\definecolor{mutedtext}    {HTML}{6B6560}  % secondary / caption text
\definecolor{rulecol}      {HTML}{C8BFB0}  % separator / rule

% Callout box palette
\definecolor{defnFrame}    {HTML}{9DB89B}  \definecolor{defnBack}    {HTML}{EEF2E8}
\definecolor{ideaFrame}    {HTML}{9FB4CD}  \definecolor{ideaBack}    {HTML}{E8EEF5}
\definecolor{warnFrame}    {HTML}{CFAE87}  \definecolor{warnBack}    {HTML}{F5EDDF}
\definecolor{noteFrame}    {HTML}{BFA6AE}  \definecolor{noteBack}    {HTML}{F2E8EE}
\definecolor{exFrame}      {HTML}{9FBDBD}  \definecolor{exBack}      {HTML}{E8F0F0}
\definecolor{codeBack}     {HTML}{F3EDE2}
\definecolor{codeBorder}   {HTML}{D4CBBA}
\definecolor{codeKeyword}  {HTML}{6B4F7A}
\definecolor{codeString}   {HTML}{8A6B4A}
\definecolor{codeComment}  {HTML}{8A857F}

% ---- Page background + base text color --------------------------------
\usepackage{pagecolor}
\pagecolor{pagebg}
\color{bodyfg}

% ---- Fonts ------------------------------------------------------------
\usepackage[T1]{fontenc}
\IfFileExists{libertine.sty}{%
    \usepackage{libertine}      % warm, readable serif
    \usepackage[scaled=0.95]{inconsolata}  % softer monospace
}{}

% ---- Hyperref ---------------------------------------------------------
\usepackage{hyperref}
\hypersetup{
    colorlinks=true,
    linkcolor=linkcol,
    urlcolor=linkcol,
    citecolor=linkcol,
    pdfborder={0 0 0},
}

% ---- Section styling --------------------------------------------------
\usepackage[explicit]{titlesec}
\titleformat{\section}
    {\color{heading}\Large\bfseries}
    {\color{heading}\thesection\hspace{0.6em}}{0pt}{#1}
\titleformat{\subsection}
    {\color{subheading}\large\bfseries}
    {\color{subheading}\thesubsection\hspace{0.5em}}{0pt}{#1}
\titleformat{\subsubsection}
    {\color{accent}\normalsize\bfseries}
    {\color{accent}\thesubsubsection\hspace{0.5em}}{0pt}{#1}
\titleformat{\paragraph}[runin]
    {\color{accent}\bfseries}{}{0pt}{#1.}

\titlespacing*{\section}{0pt}{1.6em}{0.6em}
\titlespacing*{\subsection}{0pt}{1.2em}{0.4em}

% ---- Enumerate / itemize spacing --------------------------------------
\usepackage{enumitem}
\setlist[itemize]{leftmargin=1.4em, itemsep=2pt, topsep=4pt}
\setlist[enumerate]{leftmargin=1.6em, itemsep=2pt, topsep=4pt}

% ---- Code listings ----------------------------------------------------
\usepackage{listings}
\lstdefinestyle{notescode}{
    backgroundcolor=\color{codeBack},
    basicstyle=\ttfamily\small\color{bodyfg},
    keywordstyle=\color{codeKeyword}\bfseries,
    stringstyle=\color{codeString},
    commentstyle=\color{codeComment}\itshape,
    frame=single,
    framesep=6pt,
    rulecolor=\color{codeBorder},
    showstringspaces=false,
    breaklines=true,
    xleftmargin=10pt,
    xrightmargin=10pt,
    aboveskip=8pt,
    belowskip=8pt,
}
\lstset{style=notescode}

% ---- Callout boxes (tcolorbox) ----------------------------------------
\usepackage[most]{tcolorbox}

\newtcolorbox{defnbox}[1][]{
    enhanced, breakable,
    colback=defnBack, colframe=defnFrame, coltitle=bodyfg,
    title={\bfseries Definition \textemdash\ #1},
    fonttitle=\bfseries,
    boxrule=0.5pt, arc=2pt, left=8pt, right=8pt, top=6pt, bottom=6pt,
    attach boxed title to top left={xshift=10pt, yshift=-6pt},
    boxed title style={colback=defnFrame, colframe=defnFrame, arc=2pt},
    before skip=8pt, after skip=8pt,
}
\newtcolorbox{ideabox}[1][]{
    enhanced, breakable,
    colback=ideaBack, colframe=ideaFrame, coltitle=bodyfg,
    title={\bfseries Key idea #1},
    boxrule=0.5pt, arc=2pt, left=8pt, right=8pt, top=6pt, bottom=6pt,
    attach boxed title to top left={xshift=10pt, yshift=-6pt},
    boxed title style={colback=ideaFrame, colframe=ideaFrame, arc=2pt},
    before skip=8pt, after skip=8pt,
}
\newtcolorbox{warnbox}[1][]{
    enhanced, breakable,
    colback=warnBack, colframe=warnFrame, coltitle=bodyfg,
    title={\bfseries Gotcha #1},
    boxrule=0.5pt, arc=2pt, left=8pt, right=8pt, top=6pt, bottom=6pt,
    attach boxed title to top left={xshift=10pt, yshift=-6pt},
    boxed title style={colback=warnFrame, colframe=warnFrame, arc=2pt},
    before skip=8pt, after skip=8pt,
}
\newtcolorbox{notebox}[1][]{
    enhanced, breakable,
    colback=noteBack, colframe=noteFrame, coltitle=bodyfg,
    title={\bfseries Aside #1},
    boxrule=0.5pt, arc=2pt, left=8pt, right=8pt, top=6pt, bottom=6pt,
    attach boxed title to top left={xshift=10pt, yshift=-6pt},
    boxed title style={colback=noteFrame, colframe=noteFrame, arc=2pt},
    before skip=8pt, after skip=8pt,
}
\newtcolorbox{examplebox}[1][]{
    enhanced, breakable,
    colback=exBack, colframe=exFrame, coltitle=bodyfg,
    title={\bfseries Example #1},
    boxrule=0.5pt, arc=2pt, left=8pt, right=8pt, top=6pt, bottom=6pt,
    attach boxed title to top left={xshift=10pt, yshift=-6pt},
    boxed title style={colback=exFrame, colframe=exFrame, arc=2pt},
    before skip=8pt, after skip=8pt,
}

% ---- TikZ graph styles (added 2026-04-25 for ch_10 augmentation) ------
% tcolorbox[most] already loads tikz; this block declares the libraries
% + shared `vertex` / `edge` styles used by ch_10's general directed-graph
% diagrams. ch_9's tree-style TikZ uses inline overrides and is
% unaffected. Simple circle nodes + arrow edges only -- no production
% styling.
\usetikzlibrary{arrows.meta, positioning, calc}
\tikzset{
    vertex/.style={circle, draw, minimum size=7mm, inner sep=0pt,
                   font=\small, thick},
    vsource/.style={vertex, fill=ideaBack},
    vdone/.style={vertex, fill=defnBack},
    edge/.style={->, >=Stealth, thick},
    uedge/.style={-, thick},
    treeedge/.style={->, >=Stealth, thick},
    backedge/.style={->, >=Stealth, thick, dashed, red!70!black},
    fwdedge/.style={->, >=Stealth, thick, blue!60!black},
    crossedge/.style={->, >=Stealth, thick, green!50!black, dotted},
    mstedge/.style={-, very thick, accent},
}

% ---- Title block ------------------------------------------------------
\usepackage{titling}
\pretitle{\begin{center}\color{heading}\LARGE\bfseries}
\posttitle{\end{center}\vskip -0.4em}
\preauthor{\begin{center}\color{mutedtext}\normalsize}
\postauthor{\end{center}}
\predate{\begin{center}\color{mutedtext}\small}
\postdate{\end{center}\vspace{0.4em}{\color{rulecol}\hrule height 0.4pt}\vspace{1.2em}}
` right after \documentclass
%   in each chapter's lectures.tex and notes.tex.
% =======================================================================

% ---- Layout -----------------------------------------------------------
\usepackage[margin=1in]{geometry}
\usepackage{parskip}        % space between paragraphs, no indent
\usepackage{microtype}      % subtle kerning/spacing improvements

% ---- Colors (muted palette, easy on light-sensitive eyes) -------------
\usepackage{xcolor}
\definecolor{pagebg}       {HTML}{FAF6F0}  % warm cream background
\definecolor{bodyfg}       {HTML}{3D3A36}  % warm dark gray body text
\definecolor{heading}      {HTML}{3D6A6B}  % muted teal
\definecolor{subheading}   {HTML}{5B7A9C}  % dusty blue
\definecolor{accent}       {HTML}{8A6B4A}  % warm tan
\definecolor{linkcol}      {HTML}{5B7A9C}  % same as subheading
\definecolor{mutedtext}    {HTML}{6B6560}  % secondary / caption text
\definecolor{rulecol}      {HTML}{C8BFB0}  % separator / rule

% Callout box palette
\definecolor{defnFrame}    {HTML}{9DB89B}  \definecolor{defnBack}    {HTML}{EEF2E8}
\definecolor{ideaFrame}    {HTML}{9FB4CD}  \definecolor{ideaBack}    {HTML}{E8EEF5}
\definecolor{warnFrame}    {HTML}{CFAE87}  \definecolor{warnBack}    {HTML}{F5EDDF}
\definecolor{noteFrame}    {HTML}{BFA6AE}  \definecolor{noteBack}    {HTML}{F2E8EE}
\definecolor{exFrame}      {HTML}{9FBDBD}  \definecolor{exBack}      {HTML}{E8F0F0}
\definecolor{codeBack}     {HTML}{F3EDE2}
\definecolor{codeBorder}   {HTML}{D4CBBA}
\definecolor{codeKeyword}  {HTML}{6B4F7A}
\definecolor{codeString}   {HTML}{8A6B4A}
\definecolor{codeComment}  {HTML}{8A857F}

% ---- Page background + base text color --------------------------------
\usepackage{pagecolor}
\pagecolor{pagebg}
\color{bodyfg}

% ---- Fonts ------------------------------------------------------------
\usepackage[T1]{fontenc}
\IfFileExists{libertine.sty}{%
    \usepackage{libertine}      % warm, readable serif
    \usepackage[scaled=0.95]{inconsolata}  % softer monospace
}{}

% ---- Hyperref ---------------------------------------------------------
\usepackage{hyperref}
\hypersetup{
    colorlinks=true,
    linkcolor=linkcol,
    urlcolor=linkcol,
    citecolor=linkcol,
    pdfborder={0 0 0},
}

% ---- Section styling --------------------------------------------------
\usepackage[explicit]{titlesec}
\titleformat{\section}
    {\color{heading}\Large\bfseries}
    {\color{heading}\thesection\hspace{0.6em}}{0pt}{#1}
\titleformat{\subsection}
    {\color{subheading}\large\bfseries}
    {\color{subheading}\thesubsection\hspace{0.5em}}{0pt}{#1}
\titleformat{\subsubsection}
    {\color{accent}\normalsize\bfseries}
    {\color{accent}\thesubsubsection\hspace{0.5em}}{0pt}{#1}
\titleformat{\paragraph}[runin]
    {\color{accent}\bfseries}{}{0pt}{#1.}

\titlespacing*{\section}{0pt}{1.6em}{0.6em}
\titlespacing*{\subsection}{0pt}{1.2em}{0.4em}

% ---- Enumerate / itemize spacing --------------------------------------
\usepackage{enumitem}
\setlist[itemize]{leftmargin=1.4em, itemsep=2pt, topsep=4pt}
\setlist[enumerate]{leftmargin=1.6em, itemsep=2pt, topsep=4pt}

% ---- Code listings ----------------------------------------------------
\usepackage{listings}
\lstdefinestyle{notescode}{
    backgroundcolor=\color{codeBack},
    basicstyle=\ttfamily\small\color{bodyfg},
    keywordstyle=\color{codeKeyword}\bfseries,
    stringstyle=\color{codeString},
    commentstyle=\color{codeComment}\itshape,
    frame=single,
    framesep=6pt,
    rulecolor=\color{codeBorder},
    showstringspaces=false,
    breaklines=true,
    xleftmargin=10pt,
    xrightmargin=10pt,
    aboveskip=8pt,
    belowskip=8pt,
}
\lstset{style=notescode}

% ---- Callout boxes (tcolorbox) ----------------------------------------
\usepackage[most]{tcolorbox}

\newtcolorbox{defnbox}[1][]{
    enhanced, breakable,
    colback=defnBack, colframe=defnFrame, coltitle=bodyfg,
    title={\bfseries Definition \textemdash\ #1},
    fonttitle=\bfseries,
    boxrule=0.5pt, arc=2pt, left=8pt, right=8pt, top=6pt, bottom=6pt,
    attach boxed title to top left={xshift=10pt, yshift=-6pt},
    boxed title style={colback=defnFrame, colframe=defnFrame, arc=2pt},
    before skip=8pt, after skip=8pt,
}
\newtcolorbox{ideabox}[1][]{
    enhanced, breakable,
    colback=ideaBack, colframe=ideaFrame, coltitle=bodyfg,
    title={\bfseries Key idea #1},
    boxrule=0.5pt, arc=2pt, left=8pt, right=8pt, top=6pt, bottom=6pt,
    attach boxed title to top left={xshift=10pt, yshift=-6pt},
    boxed title style={colback=ideaFrame, colframe=ideaFrame, arc=2pt},
    before skip=8pt, after skip=8pt,
}
\newtcolorbox{warnbox}[1][]{
    enhanced, breakable,
    colback=warnBack, colframe=warnFrame, coltitle=bodyfg,
    title={\bfseries Gotcha #1},
    boxrule=0.5pt, arc=2pt, left=8pt, right=8pt, top=6pt, bottom=6pt,
    attach boxed title to top left={xshift=10pt, yshift=-6pt},
    boxed title style={colback=warnFrame, colframe=warnFrame, arc=2pt},
    before skip=8pt, after skip=8pt,
}
\newtcolorbox{notebox}[1][]{
    enhanced, breakable,
    colback=noteBack, colframe=noteFrame, coltitle=bodyfg,
    title={\bfseries Aside #1},
    boxrule=0.5pt, arc=2pt, left=8pt, right=8pt, top=6pt, bottom=6pt,
    attach boxed title to top left={xshift=10pt, yshift=-6pt},
    boxed title style={colback=noteFrame, colframe=noteFrame, arc=2pt},
    before skip=8pt, after skip=8pt,
}
\newtcolorbox{examplebox}[1][]{
    enhanced, breakable,
    colback=exBack, colframe=exFrame, coltitle=bodyfg,
    title={\bfseries Example #1},
    boxrule=0.5pt, arc=2pt, left=8pt, right=8pt, top=6pt, bottom=6pt,
    attach boxed title to top left={xshift=10pt, yshift=-6pt},
    boxed title style={colback=exFrame, colframe=exFrame, arc=2pt},
    before skip=8pt, after skip=8pt,
}

% ---- TikZ graph styles (added 2026-04-25 for ch_10 augmentation) ------
% tcolorbox[most] already loads tikz; this block declares the libraries
% + shared `vertex` / `edge` styles used by ch_10's general directed-graph
% diagrams. ch_9's tree-style TikZ uses inline overrides and is
% unaffected. Simple circle nodes + arrow edges only -- no production
% styling.
\usetikzlibrary{arrows.meta, positioning, calc}
\tikzset{
    vertex/.style={circle, draw, minimum size=7mm, inner sep=0pt,
                   font=\small, thick},
    vsource/.style={vertex, fill=ideaBack},
    vdone/.style={vertex, fill=defnBack},
    edge/.style={->, >=Stealth, thick},
    uedge/.style={-, thick},
    treeedge/.style={->, >=Stealth, thick},
    backedge/.style={->, >=Stealth, thick, dashed, red!70!black},
    fwdedge/.style={->, >=Stealth, thick, blue!60!black},
    crossedge/.style={->, >=Stealth, thick, green!50!black, dotted},
    mstedge/.style={-, very thick, accent},
}

% ---- Title block ------------------------------------------------------
\usepackage{titling}
\pretitle{\begin{center}\color{heading}\LARGE\bfseries}
\posttitle{\end{center}\vskip -0.4em}
\preauthor{\begin{center}\color{mutedtext}\normalsize}
\postauthor{\end{center}}
\predate{\begin{center}\color{mutedtext}\small}
\postdate{\end{center}\vspace{0.4em}{\color{rulecol}\hrule height 0.4pt}\vspace{1.2em}}


\title{CS 300 -- Chapter 3 Lectures\\\large Introduction to Data Structures}
\author{Jose de Lima}
\date{\today}

\begin{document}
\maketitle

\begin{ideabox}[Chapter map]
\textbf{Where this sits in CS 300.} Chapter 3 is the longest chapter in the
course and the densest conceptually. It does four things in sequence:
(1) grounds data representation (bits, ASCII) so the later structures
have a floor; (2) introduces \textbf{data structures} as named shapes for
storing data; (3) introduces \textbf{abstract data types} as the
interface-level view; (4) builds the vocabulary of \textbf{algorithm
runtime}: linear vs.\ binary search, Big-O, and the sorting zoo. Every
later chapter in this course assumes you can read ``$O(n \log n)$'' and
understand \emph{why}.

\textbf{What you add to your toolkit here:}
\begin{itemize}
    \item \textbf{A map of the course.} The seven structures (array,
          linked list, tree, BST, heap, hash table, graph) are named and
          contrasted before you build any of them.
    \item \textbf{ADT thinking.} Stacks, queues, sets, maps -- the
          contract (interface) is what users care about; the data
          structure (implementation) is what you choose.
    \item \textbf{Big-O fluency.} You can read, write, and justify
          complexity claims: linear search is $O(n)$ because it touches
          every element; binary search is $O(\log n)$ because each step
          halves the remaining range.
    \item \textbf{Seven sorting algorithms} with their tradeoffs.
          Insertion, selection, shell (incremental); quicksort, merge
          sort (divide-and-conquer); radix (non-comparison). Knowing
          \emph{which} to reach for is a high-frequency exam question.
\end{itemize}

\textbf{Mastery checklist} -- you should be able to, without references:
\begin{enumerate}
    \item Translate between bits, hex, and an ASCII character (e.g.,
          \texttt{0x41} $\to$ \texttt{'A'}).
    \item State the interface of each ADT (List, Stack, Queue, Set, Map,
          Priority Queue) and give one concrete data structure that
          implements it.
    \item Explain linear vs.\ binary search with the runtime for each on
          the worst case and the preconditions.
    \item State the Big-O definition from memory (``$f(n) = O(g(n))$
          iff there exist constants $c$ and $n_0 \ldots$'') and apply it
          to a small example.
    \item Order the growth hierarchy: $1, \log n, n, n \log n, n^2, 2^n$.
    \item For each sort (insertion, selection, shell, quicksort, merge,
          radix), give best / average / worst / stable? / in-place?
          \emph{and} the one sentence that explains the runtime.
    \item Pick the right sort for a scenario: ``small, nearly sorted,''
          ``huge, RAM-bound,'' ``keys are integers with bounded range,''
          ``must be stable.''
\end{enumerate}

\textbf{Looking ahead.} Chapters 4-6 spend their entire budget building
two or three of the structures from the list here (linked lists, hash
tables, BSTs). Every performance claim in those chapters is phrased in
Big-O, and most of the sort algorithms in 3.15-3.21 reappear as
components (e.g., merge sort is the model for merging sorted linked
lists in ch.~4).
\end{ideabox}

\section{3.1 Representing Information as Bits}

Chapter 2 covered algorithms -- the \emph{instructions} that process
data. Chapter 3 shifts the view to how data is \emph{represented} in
the first place. At the bottom of every stack, ``data'' is ultimately
a sequence of tiny binary switches. This section is the
ground-floor view: what a bit is, what a byte is, and how characters
like \texttt{'A'} or \texttt{'Z'} get mapped onto patterns of bits.

\subsection{Bits and bytes}

\begin{defnbox}[Bit]
A \textbf{bit} (``binary digit'') is the smallest unit of information:
a single \texttt{0} or \texttt{1}. Physically, it corresponds to a
switch, a charged capacitor, a magnetic polarity, or a voltage level
-- anything with two stable states.
\end{defnbox}

\begin{defnbox}[Byte]
A \textbf{byte} is a group of \textbf{8 bits}, giving $2^8 = 256$
distinct values (bit patterns \texttt{00000000} through
\texttt{11111111}). Memory, file sizes, and network transfers are
almost always measured in bytes or larger units (KB, MB, GB \ldots).
\end{defnbox}

\begin{notebox}[Nibbles, halfwords, and other groupings]
A \textbf{nibble} is 4 bits (half a byte). A \textbf{word} is typically
what the CPU processes in one step -- 4 bytes (32-bit) or 8 bytes
(64-bit) on modern machines. C++ types you've already used map onto
these sizes: \texttt{char} is 1 byte, \texttt{int} is usually 4, a
pointer or \texttt{size\_t} on a 64-bit machine is 8. \texttt{sizeof(T)}
tells you the number of bytes.
\end{notebox}

\subsection{Why bits}

Every piece of data -- an integer, a character, a pixel, a video frame,
a neural network's weights -- gets ultimately stored as a pattern of
bits. The question isn't \emph{whether} data is bits but \emph{which
encoding} the program uses to interpret those bits. Two different
encodings can give the same pattern two different meanings:

\begin{examplebox}[\texttt{01000001} -- what does it mean?]
\begin{itemize}[leftmargin=*,noitemsep]
    \item As an unsigned 8-bit integer: $65$.
    \item As a signed 8-bit integer (two's complement): $65$.
    \item As an ASCII character: \texttt{'A'}.
    \item As a pair of nibbles: $4$ and $1$.
    \item As a bitfield: bit 6 on, bits 0 and 5 on, rest off.
\end{itemize}
Same eight bits, five different interpretations. The program decides
which one applies based on the type declaration.
\end{examplebox}

\subsection{ASCII: bits that mean characters}

\begin{defnbox}[Character encoding]
A \textbf{character encoding} is a table mapping numeric codes to
characters (letters, digits, punctuation, symbols). To ``store a
string,'' the computer stores the sequence of code values; when it
``prints a string,'' it looks up the characters those codes represent.
\end{defnbox}

\begin{defnbox}[ASCII]
\textbf{ASCII} (\emph{American Standard Code for Information
Interchange}) is the oldest widely-used character encoding. It
assigns a unique 7-bit code (values \texttt{0}--\texttt{127}) to each
printable English character plus a set of control codes. Because it
fits in 7 bits, every ASCII character fits comfortably in a single
byte with one bit to spare.
\end{defnbox}

\subsection{A subset of the ASCII table}

The full ASCII table is 128 entries. A few useful patterns to commit
to memory:

\begin{center}
\small
\begin{tabular}{lll}
\textbf{Range (dec)} & \textbf{Range (hex)} & \textbf{Meaning} \\
\hline
0--31  & \texttt{0x00}--\texttt{0x1F} & Control characters (non-printing) \\
32     & \texttt{0x20}                & Space \\
33--47 & \texttt{0x21}--\texttt{0x2F} & Punctuation: \texttt{! " \# \$ \% \& ' ( ) * + , - . /} \\
48--57 & \texttt{0x30}--\texttt{0x39} & Digits \texttt{'0'}--\texttt{'9'} \\
58--64 & \texttt{0x3A}--\texttt{0x40} & More punctuation: \texttt{: ; < = > ? @} \\
65--90 & \texttt{0x41}--\texttt{0x5A} & Uppercase \texttt{'A'}--\texttt{'Z'} \\
91--96 & \texttt{0x5B}--\texttt{0x60} & \texttt{[ \textbackslash\ ] \textasciicircum\ \_ `} \\
97--122 & \texttt{0x61}--\texttt{0x7A} & Lowercase \texttt{'a'}--\texttt{'z'} \\
123--126 & \texttt{0x7B}--\texttt{0x7E} & \texttt{\{ | \} \textasciitilde} \\
127 & \texttt{0x7F} & DEL (control) \\
\end{tabular}
\normalsize
\end{center}

\begin{ideabox}[The three patterns worth remembering]
\begin{enumerate}[leftmargin=*,noitemsep]
    \item \texttt{'0'} is $48$, \texttt{'1'} is $49$, \ldots, \texttt{'9'} is $57$.
          \textbf{\texttt{c - '0'}} converts a digit character to its integer value.
    \item \texttt{'A'} is $65$, \texttt{'Z'} is $90$; \texttt{'a'} is
          $97$, \texttt{'z'} is $122$. Uppercase and lowercase differ
          by exactly \textbf{32} ($2^5$), so \texttt{c + 32} lowercases an
          uppercase letter and \texttt{c - 32} uppercases a lowercase
          one. \texttt{c \textasciicircum\ 32} does both (toggle case).
    \item Space is $32$ = $0x20$. Every printable character is $\ge 32$.
          Anything below is a control character (newline is $10$,
          carriage return is $13$, tab is $9$, null is $0$).
\end{enumerate}
These three facts alone cover half of the character-arithmetic bugs
you'll hit.
\end{ideabox}

\begin{examplebox}[Digit-char to int, by hand]
\begin{verbatim}
char c = '7';
int  d = c - '0';     // 55 - 48 == 7

// Accumulate digits from a string into a number:
int n = 0;
for (char ch : s) {
    if (std::isdigit(ch)) n = n * 10 + (ch - '0');
}
\end{verbatim}
\end{examplebox}

\begin{examplebox}[Case conversion without \texttt{cctype}]
\begin{verbatim}
char upper = 'g' - 32;     // 'G'
char lower = 'G' + 32;     // 'g'
char flip  = 'G' ^ 32;     // 'g'
char flip2 = 'g' ^ 32;     // 'G'
\end{verbatim}
Using \texttt{std::toupper}/\texttt{std::tolower} is still the right
production choice -- they handle the non-letter case and locale
concerns correctly -- but knowing the underlying bit trick helps when
you read code that uses it.
\end{examplebox}

\subsection{What ASCII can't do: everything else}

ASCII encodes English. It has no code points for accented letters
(\texttt{é}, \texttt{ñ}), no space for non-Latin scripts (Arabic,
Chinese, Cyrillic, Tamil), and no emoji. The modern universal encoding
is \textbf{Unicode}, which assigns code points up to $10FFFF_{16}$
($\sim 1.1$ million characters).

\begin{defnbox}[Unicode and UTF-8]
\textbf{Unicode} is the catalog of characters and their numeric code
points. \textbf{UTF-8} is the most common byte-level encoding: ASCII
characters still fit in one byte (and use the same values as ASCII!),
while higher code points use 2--4 bytes. UTF-8 is backward compatible
with ASCII, which is why almost every file, web page, and API
defaults to it.
\end{defnbox}

\begin{warnbox}[\texttt{char} vs.\ Unicode in C++]
A \texttt{char} in C++ is one byte. A UTF-8 character (like
\texttt{é}) can occupy multiple bytes. Code that walks a UTF-8 string
one \texttt{char} at a time is counting bytes, not characters, and
treating each byte as an independent character will produce garbage
output and wrong indices for any non-ASCII input. C++ offers
\texttt{char16\_t}, \texttt{char32\_t}, and \texttt{std::u8string} for
wider character work, plus libraries like ICU for correct text
handling.

For CS 300 projects with English-only input you rarely need to worry
-- but the assumption ``1 char = 1 character'' is only true for
ASCII, and ``is only true for ASCII'' means it's only true for one
small slice of the world's text.
\end{warnbox}

\subsection{Connecting bits back to data structures}

Every data structure you're about to meet (linked list, stack, queue,
tree, hash table) is ultimately a layout of bits in memory, organized
so that the operations the structure promises are fast. A linked list
stores a value plus a pointer-to-next-node; a vector stores an array
of values plus a size and capacity; a hash table stores an array of
buckets plus a hash function. The abstractions hide the bits, but
understanding that the bits are still there -- and that performance
ultimately comes from how they're arranged -- is what separates
``using'' a data structure from \emph{choosing} one.

\begin{ideabox}[Big picture]
All data is bits; all bit patterns need an encoding to become data.
ASCII is the oldest and simplest encoding, covering exactly the
characters on a 1960s American typewriter, and is still the
foundation of modern text encodings that extend it to the rest of
human writing. The arithmetic shortcuts -- \texttt{c - '0'} for
digits, \texttt{\textasciicircum\ 32} for case, \texttt{c >= '0' \&\&
c <= '9'} for ``is digit'' -- come from ASCII's tidy layout and show
up constantly in real parsing code. Build intuition for these now;
every chapter from here on assumes it.
\end{ideabox}

\section{3.2 Data Structures}

A \textbf{data structure} is a way of organizing data in memory so
that a particular set of operations (access, search, insert, delete)
can be done efficiently. \textsection 3.1 looked at how individual
values are stored as bits; this section starts naming the ways those
values get grouped into useful collections. The rest of the course is
essentially a detailed tour of the seven structures listed here.

\begin{defnbox}[Data structure]
A \textbf{data structure} is a combination of a \emph{layout in
memory} and a \emph{set of operations} that together give certain
common tasks good performance characteristics. The same data can be
organized many ways -- the choice depends on which operations matter
most for your workload.
\end{defnbox}

\subsection{The seven structures you'll meet this semester}

\begin{defnbox}[Record (struct, class)]
A group of named \textbf{fields} bundled into a single value. Fields
can be of different types. In C++, a \texttt{struct} or \texttt{class}
is a record.
\begin{verbatim}
struct Bid {
    std::string title;
    std::string fund;
    double amount;
};
\end{verbatim}
\end{defnbox}

\begin{defnbox}[Array / vector]
An \textbf{ordered} list of items stored in \emph{contiguous memory},
directly indexable by position. Access is $O(1)$; insert/delete in
the middle is $O(n)$ because later items have to shift.
\end{defnbox}

\begin{defnbox}[Linked list]
An ordered list stored as a chain of \textbf{nodes}, each holding a
value and a pointer to the next node. No contiguous memory, no
indexing. Insert/delete given a pointer to the node is $O(1)$;
accessing the $k$-th item is $O(k)$ because you have to walk the
chain.
\end{defnbox}

\begin{defnbox}[Binary tree]
A hierarchical structure where each node has at most two children
(left and right). The foundation for binary search trees, heaps,
expression trees, and more. Good implementations give $O(\log n)$
search, insert, and delete on balanced variants.
\end{defnbox}

\begin{defnbox}[Hash table]
An array plus a \textbf{hash function} that maps keys to array
positions. Average-case $O(1)$ for insert, lookup, and delete; no
ordering. The standard choice when you need fast keyed access and
don't need sorted iteration.
\end{defnbox}

\begin{defnbox}[Heap]
A tree maintaining the \textbf{heap property}: in a \emph{max-heap},
each node's key is $\ge$ its children's keys; in a \emph{min-heap},
$\le$. Gives $O(\log n)$ insert and remove-extreme. The underlying
structure behind priority queues and heapsort.
\end{defnbox}

\begin{defnbox}[Graph]
A collection of \textbf{vertices} connected by \textbf{edges}, used
to model networks (roads, social graphs, dependencies, web links).
Can be directed or undirected, weighted or unweighted. Operations
include traversal, shortest path, and connectivity.
\end{defnbox}

\subsection{Why there are so many}

Every structure makes different operations fast at the cost of others.
There is no ``best'' data structure; the right one depends on the
operation mix you actually perform.

\begin{center}
\small
\begin{tabular}{p{3.3cm} p{2.4cm} p{2.4cm} p{3.6cm}}
\textbf{Structure} & \textbf{Access} & \textbf{Insert mid} &
\textbf{Good for} \\
\hline
Array / vector & $O(1)$ & $O(n)$ & Random access, cache-friendly \\
Linked list & $O(n)$ & $O(1)$\textsuperscript{*} & Frequent insert/delete \\
Hash table & $O(1)$ avg & $O(1)$ avg & Fast keyed lookup \\
Binary search tree & $O(\log n)$\textsuperscript{\dag} &
$O(\log n)$\textsuperscript{\dag} & Sorted iteration \\
Heap & --- & $O(\log n)$ & Always know the max/min \\
Graph & --- & varies & Relationships, networks \\
\end{tabular}
\end{center}
\small
\textsuperscript{*}Given a pointer to the insertion point. Finding the
point is still $O(n)$.

\textsuperscript{\dag}Balanced only. A degenerate BST is $O(n)$.
\normalsize

\subsection{A concrete example: the shift problem}

The classic case where structure choice matters is inserting in the
middle of a collection.

\begin{examplebox}[Array insertion is $O(n)$]
Start: \texttt{[A, C, D]}. Insert \texttt{B} at position 1.

To make room, \texttt{C} and \texttt{D} have to shift right:
\texttt{[A, \_, C, D]} $\to$ \texttt{[A, \_, \_, C, D]} $\to$
\texttt{[A, B, C, D]}.

Every item after the insertion point moves. For an array of size $n$,
worst-case insertion moves $n$ items.
\end{examplebox}

\begin{examplebox}[Linked list insertion is $O(1)$ (given the node)]
Start: \texttt{A $\to$ C $\to$ D}. To insert \texttt{B} between
\texttt{A} and \texttt{C}:

\begin{enumerate}[leftmargin=*,noitemsep]
    \item Allocate a new node \texttt{B}.
    \item Point \texttt{B.next} at the node \texttt{C}.
    \item Point \texttt{A.next} at \texttt{B}.
\end{enumerate}
Two pointer writes. No shifting, no matter how long the list is.
\end{examplebox}

\begin{warnbox}[``$O(1)$ linked list insert'' has a hidden assumption]
The $O(1)$ cost assumes you already have a pointer to the node where
you want to insert. If you have to \emph{find} that node first --
``insert B after the node containing A'' -- you pay $O(n)$ to find
it, then $O(1)$ to do the insert. For insertion at an arbitrary
position by value, linked lists usually don't beat vectors in
practice.
\end{warnbox}

\subsection{How to pick a data structure}

\begin{ideabox}[A practical decision procedure]
\begin{enumerate}[leftmargin=*,noitemsep]
    \item \textbf{What operations dominate?} Access by index, by key,
          by position? Insert at the ends, or in the middle? Iterate
          in sorted order? Always need the min/max?
    \item \textbf{What's the expected size?} Small ($\le 100$) almost
          anything is fine and a vector usually wins by cache alone.
          Large means the asymptotic cost matters.
    \item \textbf{What invariant do you need?} Sorted? Unique? Order
          of insertion preserved? That usually rules out several
          structures.
    \item \textbf{Is the standard-library version good enough?}
          \texttt{std::vector}, \texttt{std::unordered\_map},
          \texttt{std::set}, and \texttt{std::priority\_queue} cover
          most needs -- reach for a custom structure only when
          measurement shows the defaults are the bottleneck.
\end{enumerate}
\end{ideabox}

\begin{examplebox}[Mapping three typical CS 300 tasks to structures]
\begin{itemize}[leftmargin=*,noitemsep]
    \item ``Look up a course by its ID.'' $\to$ hash table
          (\texttt{std::unordered\_map<std::string, Course>}).
    \item ``Print courses in alphabetical order.'' $\to$ balanced BST
          (\texttt{std::map}) or a vector + sort once at the end.
    \item ``Process bids in arrival order.'' $\to$ queue (backed by a
          linked list or a ring buffer).
\end{itemize}
Same data, different access pattern, different structure.
\end{examplebox}

\subsection{The C++ standard library's version of this list}

Almost every structure named above already ships with the language.
Know the mapping so you don't reinvent:

\begin{center}
\small
\begin{tabular}{ll}
\textbf{Concept} & \textbf{C++ standard type} \\
\hline
Array / vector & \texttt{std::vector<T>}, \texttt{std::array<T, N>} \\
Linked list & \texttt{std::list<T>} (doubly) or \texttt{std::forward\_list<T>} \\
Hash table & \texttt{std::unordered\_map<K, V>}, \texttt{std::unordered\_set<T>} \\
Balanced BST & \texttt{std::map<K, V>}, \texttt{std::set<T>} \\
Heap / priority queue & \texttt{std::priority\_queue<T>} \\
Stack / queue (adapters) & \texttt{std::stack<T>}, \texttt{std::queue<T>} \\
Graph & no built-in -- build from \texttt{vector<vector<int>>} or similar \\
\end{tabular}
\normalsize
\end{center}

\begin{notebox}[Build your own for understanding, use the STL in practice]
CS 300 will have you implement versions of several of these from
scratch. That's for learning. In real code, reach for the library
first; it's tuned, tested, and portable. Understanding what's
happening inside \texttt{std::unordered\_map} makes you better at
using it, not a replacement for using it.
\end{notebox}

\begin{ideabox}[Big picture]
The rest of the course is a depth-first walk through these seven
structures -- their implementations, their asymptotic costs, and the
patterns of problem each one fits. Think of this section as the map:
every later chapter picks one structure, builds it in detail,
analyzes its operations, and shows the problems it solves. When you
finish, the skill is not ``I've memorized all seven'' but ``I can
look at a problem and say, that's a bag-of-key-value pairs with fast
lookup -- hash table.''
\end{ideabox}

\section{3.3 Relation Between Data Structures and Algorithms}

Chapter 2 treated algorithms in the abstract; section 3.2 laid out the
main data structures. This section is the hinge between them. The
relationship runs in \emph{two directions}, and keeping both in mind
is what turns ``I know some algorithms'' and ``I know some
structures'' into ``I can design software.''

\begin{ideabox}[Two directions of the same relationship]
\begin{enumerate}[leftmargin=*]
  \item \textbf{Algorithms \emph{for} a data structure.} Each structure
        needs its own implementation of common operations (insert,
        remove, search). The \emph{same operation name} hides very
        different algorithms depending on the underlying structure.
  \item \textbf{Algorithms that \emph{use} data structures.} A
        larger algorithm typically leans on one or more structures as
        scratch space -- somewhere to accumulate, sort, or look up
        data while the algorithm runs.
\end{enumerate}
\end{ideabox}

\subsection{Direction 1: Different structure, different algorithm}

``Append an item'' sounds like one operation, but the code behind it
depends entirely on what you're appending to. Compare appending to a
linked list against appending to an array.

\begin{examplebox}[Append to a singly linked list -- \texttt{O(1)}]
\begin{lstlisting}[basicstyle=\ttfamily\small,frame=single,language=C++]
void append(Node*& head, Node*& tail, int value) {
    Node* n = new Node{value, nullptr};
    if (!head) {           // empty list
        head = tail = n;
    } else {
        tail->next = n;    // stitch onto the end
        tail = n;
    }
}
\end{lstlisting}
With a tail pointer cached, linking on a new node is a couple of
pointer writes -- no elements moved, no memory copied.
\end{examplebox}

\begin{examplebox}[Append to a fixed-capacity array -- \texttt{O(1)} amortized, \texttt{O(n)} worst case]
\begin{lstlisting}[basicstyle=\ttfamily\small,frame=single,language=C++]
void append(int*& data, int& size, int& cap, int value) {
    if (size == cap) {
        int newCap = cap == 0 ? 1 : cap * 2;
        int* grown = new int[newCap];
        for (int i = 0; i < size; ++i) grown[i] = data[i];  // copy all
        delete[] data;
        data = grown;
        cap = newCap;
    }
    data[size++] = value;
}
\end{lstlisting}
Most appends just drop the value into the next slot. But whenever the
array fills up, every existing element has to be copied into a larger
block. \texttt{std::vector} hides this dance behind \texttt{push\_back},
but the copy still happens.
\end{examplebox}

\begin{notebox}[``Amortized'' in one sentence]
If doubling on overflow makes some individual pushes expensive but
\emph{averages out} to constant time across many pushes, we say
\texttt{push\_back} is \textbf{amortized \(O(1)\)}. You pay a big
cost occasionally, but the cost per call, averaged over a long run,
stays constant. We'll make this precise when we cover analysis.
\end{notebox}

\subsection*{Insert-at-front is the sharper example}

Append hides a lot because both versions \emph{usually} finish fast.
Prepend makes the structure choice visible immediately.

\begin{center}
\footnotesize
\renewcommand{\arraystretch}{1.25}
\begin{tabular}{@{}lll@{}}
\hline
Operation & Array / \texttt{vector} & Singly linked list \\
\hline
Append at end        & $O(1)$ amortized    & $O(1)$ with tail pointer \\
Insert at front      & $O(n)$ (shift all)  & $O(1)$ (new head) \\
Random access \texttt{a[i]} & $O(1)$         & $O(i)$ (walk from head) \\
Insert in middle (given iterator) & $O(n)$ shift & $O(1)$ relink \\
\hline
\end{tabular}
\end{center}

Same \emph{operation names}, completely different runtimes. The
structure you pick \emph{is} the set of costs you accept.

\begin{warnbox}[``Same name'' is misleading]
Don't trust operation names across containers. \texttt{push\_back}
on a \texttt{std::vector} is fast; \texttt{push\_front} doesn't even
exist (use \texttt{std::deque} or \texttt{std::list}).
\texttt{find} on \texttt{std::vector} is $O(n)$; on
\texttt{std::unordered\_map} it's expected $O(1)$. Always ask: what
is the cost of this operation \emph{on this structure}?
\end{warnbox}

\subsection{Direction 2: Algorithms use structures as scratch space}

The other direction: a larger algorithm almost always relies on a data
structure to hold intermediate results. The structure you pick shapes
the algorithm you write.

\textbf{Example: top-\(k\) salespersons.} Given a list of salespeople,
print the top 5 by sales total.

\begin{examplebox}[A small fixed-size array as running top-5]
\begin{lstlisting}[basicstyle=\ttfamily\small,frame=single,language=C++]
struct Rep { std::string name; double sales; };

void topFive(const std::vector<Rep>& all) {
    std::array<Rep, 5> top{};
    for (auto& r : top) r.sales = -1;        // sentinel

    for (const Rep& r : all) {
        if (r.sales > top.back().sales) {    // beats last in top-5
            top.back() = r;
            std::sort(top.begin(), top.end(),
                      [](const Rep& a, const Rep& b){
                          return a.sales > b.sales;
                      });
        }
    }
    for (const Rep& r : top) std::cout << r.name << " " << r.sales << "\n";
}
\end{lstlisting}
The fixed-size array is the data structure; the scan-and-swap pattern
is the algorithm. The array's job is to remember the current champions
as we sweep through all candidates.
\end{examplebox}

\begin{ideabox}[Why fixed-size-5 instead of ``sort them all''?]
Sorting the entire list is \(O(n \log n)\). The running-top-5 sweep is
\(O(n \log k)\) where \(k = 5\) -- effectively \(O(n)\) because the
inner sort is on a constant-size array. For top-$k$ problems where
$k \ll n$, this is the standard trick. (A heap of size $k$ is even
better; we'll see that in the priority-queue chapter.)
\end{ideabox}

\subsection*{Different structure, different algorithm}

If we used a different structure to hold the running top-5, the
algorithm changes shape:
\begin{itemize}[leftmargin=*]
  \item \textbf{Min-heap of size 5.} Peek at the minimum in $O(1)$;
        if the new candidate beats it, pop and push. Total
        $O(n \log k)$ but cleaner and faster than re-sorting an array.
  \item \textbf{Sorted linked list of size 5.} Insert-in-order means
        walking the list to find the right slot -- $O(k)$ per
        candidate that beats the minimum.
  \item \textbf{BST / \texttt{std::set}.} Overkill here, but useful if
        $k$ were large or you needed dynamic queries.
\end{itemize}

Same problem, four different implementations, each dictated by the
structure you choose.

\begin{warnbox}[Structure choice $\neq$ premature optimization]
``Pick the simplest thing that works'' is good advice. But the
structure is often what determines whether ``works'' is even feasible
at your data size. A naive $O(n^2)$ search that's \emph{correct} is
fine for 100 items and ruinous for 100\,000\,000. The choice is part
of the design, not a polish step at the end.
\end{warnbox}

\subsection{Putting the two directions together}

\begin{ideabox}[What it means to ``design'' something]
When you sit down with a new problem, the workflow is:
\begin{enumerate}[leftmargin=*]
  \item \textbf{What data does the algorithm need access to, and how?}
        By index? By key? In sorted order? First-in-first-out?
  \item \textbf{Which structure answers those access patterns cheaply?}
        Consult the table from 3.2.
  \item \textbf{Now: what does the algorithm look like given that
        structure's operations?} Its cost is the sum of the per-step
        costs in the structure you picked.
\end{enumerate}
Step 2 is the choice \emph{of} structure; step 3 is the algorithm
\emph{using} it. Those are the two directions, in order, every time.
\end{ideabox}

\begin{notebox}[Where this is going]
The rest of the course iterates this loop structure-by-structure. For
each of the seven structures from 3.2, we'll define it precisely,
implement it in C++, analyze the cost of its core operations, and
study the algorithms that either implement it or use it. By the end,
the ``design'' step above will be muscle memory.
\end{notebox}

\section{3.4 Abstract Data Types}

The previous section showed that the \emph{same operation name} on
two different data structures hides two different algorithms. That's
a feature, not a bug. The whole point of this section's central idea
-- the \textbf{abstract data type} -- is to let us \emph{name an
operation without naming an implementation}, so code that uses the
operation keeps working no matter which structure sits underneath.

\begin{defnbox}[Abstract data type (ADT)]
An \textbf{abstract data type} is a data type defined by the
\emph{operations} it supports and the \emph{behavior} of those
operations, \textbf{without} specifying how they are implemented.
An ADT is a contract: ``you can do these things, and they mean this.''
A \textbf{data structure} is one concrete way to honor that contract.
\end{defnbox}

\begin{ideabox}[ADT vs.\ data structure, in one line]
An \textbf{ADT} is \emph{what}. A \textbf{data structure} is
\emph{how}. A single ADT can have many structures behind it; a single
structure can be used to implement several ADTs.
\end{ideabox}

\subsection{Why the separation matters}

Everyday analogy: ``printer'' is an ADT (print this page, cancel this
job). Laser printer, inkjet, thermal, network printer -- all different
implementations. The program calling \texttt{print()} does not and
should not care. Swap the printer, the program still works.

In code, the benefits land as:
\begin{itemize}[leftmargin=*]
  \item \textbf{Replaceability.} Start with an array-backed list;
        switch to a linked-list-backed list when profiling tells you
        to. No caller code changes.
  \item \textbf{Reasoning.} The correctness of caller code depends on
        the ADT's \emph{contract}, not on representation details.
  \item \textbf{Teamwork.} One person builds the ADT, another builds
        against its interface. They negotiate a contract, not a
        memory layout.
\end{itemize}

\begin{warnbox}[Contract includes \emph{cost}]
A pure contract-of-behavior is not quite enough in this course. When
we promise ``\(O(1)\) lookup,'' the constant-time guarantee is part
of the ADT's contract. Two implementations that both return the right
answer can still be usefully different ADTs if their cost profiles
differ -- which is why \texttt{std::map} (balanced BST, $O(\log n)$)
and \texttt{std::unordered\_map} (hash table, expected $O(1)$) are
separate types rather than one ``dictionary'' type with a hidden
switch.
\end{warnbox}

\subsection{The standard ADTs you should recognize}

The textbook lists nine common ADTs. Treat this table as a
vocabulary list -- each one gets its own chapter later.

\begin{center}
\footnotesize
\renewcommand{\arraystretch}{1.3}
\begin{tabular}{@{}p{2.7cm}p{6.1cm}p{4.2cm}@{}}
\hline
\textbf{ADT} & \textbf{What the contract says} & \textbf{Usual implementations} \\
\hline
List          & Ordered sequence; append, remove, search, traverse.      & Array, linked list \\
Dynamic array & Ordered with \(O(1)\) indexed access; grows on demand.   & Array (with resize) \\
Stack         & LIFO; push / pop / peek at the top only.                 & Array, linked list \\
Queue         & FIFO; enqueue at back, dequeue from front.               & Linked list, circular array \\
Deque         & Insert/remove at \emph{both} ends.                       & Doubly linked list, circular array \\
Bag           & Unordered collection; duplicates allowed.                & Array, linked list \\
Set           & Unordered; \emph{no} duplicates.                         & BST, hash table \\
Priority queue & Queue where highest-priority item comes out first.     & Heap \\
Dictionary / Map & Key \(\to\) value lookup.                             & Hash table, BST \\
\hline
\end{tabular}
\end{center}

\begin{notebox}[How to read the rows]
Two axes usually disambiguate them: (1) is order required, and if so
what kind (insertion order, priority order, sorted by key)? (2) are
duplicates allowed? A \emph{bag} has duplicates and no order; a
\emph{set} has no duplicates and no order; a \emph{list} has duplicates
and insertion order; a \emph{priority queue} has duplicates and
priority order.
\end{notebox}

\subsection{C++ makes the ADT/structure distinction concrete}

The C++ standard library names ADTs via \emph{container adaptors}
that sit on top of underlying containers. This is worth seeing
explicitly because it's the same idea the textbook is describing,
except made literal in the type system.

\begin{examplebox}[\texttt{std::stack} is an ADT built on another container]
\begin{lstlisting}[basicstyle=\ttfamily\small,frame=single,language=C++]
#include <stack>
#include <vector>
#include <list>

std::stack<int>                     s1;  // default: uses std::deque
std::stack<int, std::vector<int>>   s2;  // uses a vector
std::stack<int, std::list<int>>     s3;  // uses a linked list

s1.push(3); s1.push(4);
int top = s1.top();  // 4
s1.pop();
\end{lstlisting}
All three behave identically from the caller's point of view: a
stack. What differs is the underlying storage. The ADT (\texttt{stack})
is the interface; the structure is a template parameter.
\end{examplebox}

\subsection*{The library's ADT / structure mapping}

\begin{center}
\footnotesize
\renewcommand{\arraystretch}{1.25}
\begin{tabular}{@{}llp{4.5cm}@{}}
\hline
\textbf{ADT (from textbook)} & \textbf{C++ name} & \textbf{Default structure} \\
\hline
List              & \texttt{std::list}, \texttt{std::forward\_list} & Doubly / singly linked list \\
Dynamic array     & \texttt{std::vector}           & Contiguous array \\
Stack             & \texttt{std::stack}            & \texttt{std::deque} \\
Queue             & \texttt{std::queue}            & \texttt{std::deque} \\
Deque             & \texttt{std::deque}            & Chunked array \\
Set               & \texttt{std::set}, \texttt{std::unordered\_set} & BST / hash table \\
Priority queue    & \texttt{std::priority\_queue}  & Binary heap over \texttt{std::vector} \\
Map               & \texttt{std::map}, \texttt{std::unordered\_map} & BST / hash table \\
\hline
\end{tabular}
\end{center}

(C++ doesn't ship a ``bag'' type; a \texttt{std::vector} or
\texttt{std::multiset} fills that role.)

\begin{warnbox}[Don't confuse \texttt{std::vector} and ``array'']
C++'s built-in array (\texttt{int a[10]}) is a fixed-size data
structure with no resize. \texttt{std::vector} is the ADT ``dynamic
array,'' implemented \emph{using} a C array plus a size and capacity
counter. When someone says ``array'' in conversation about C++,
they usually mean \texttt{std::vector}. Read carefully.
\end{warnbox}

\subsection{Using ADTs to keep caller code honest}

\begin{examplebox}[Algorithm expressed against an ADT]
\begin{lstlisting}[basicstyle=\ttfamily\small,frame=single,language=C++]
// Reverse the order of items in any container via a stack.
template <typename Container>
void reverseInPlace(Container& c) {
    std::stack<typename Container::value_type> buf;
    for (auto& x : c) buf.push(x);   // push all
    for (auto& x : c) {              // pop back into place
        x = buf.top();
        buf.pop();
    }
}
\end{lstlisting}
This works for any container. The algorithm depends only on the
\emph{stack} ADT -- push, top, pop. Whatever storage
\texttt{std::stack} uses, the reverse works.
\end{examplebox}

\begin{ideabox}[The ADT-first habit]
When you're sketching a solution, name the ADT you need
(``I need a priority queue'') \emph{before} you pick a structure
(``I'll back it with a binary heap''). Most design mistakes at this
level come from picking a structure first, then contorting the
algorithm around it. The fix: force yourself to say the ADT out
loud first.
\end{ideabox}

\begin{notebox}[Where this is going]
Starting in Chapter 4, every chapter for the rest of the course
follows the same pattern: pick an ADT, pick a structure (or two) that
implement it, build those implementations in C++, and analyze each
operation's cost. Sections 3.2 -- 3.4 gave you the vocabulary to
recognize which ADT/structure pairing is in front of you at any
given moment.
\end{notebox}

\section{3.5 Applications of ADTs}

Section 3.4 made the ADT/structure distinction; this section asks
what you \emph{get} from keeping them separate. Two payoffs:
\textbf{abstraction} simplifies programming, and \textbf{standard
libraries} mean you rarely implement the common ADTs by hand. The
catch -- which this section also makes explicit -- is that abstraction
does not free you from thinking about efficiency.

\subsection{Abstraction as a productivity lever}

\begin{defnbox}[Abstraction]
\textbf{Abstraction} is the practice of interacting with something
through a \emph{high-level interface}, with lower-level details
hidden behind that interface. ADTs are one concrete application of
the idea: the operations are the interface, the structure is hidden.
\end{defnbox}

Abstraction buys three compounding benefits once you're comfortable
with it:
\begin{itemize}[leftmargin=*]
  \item \textbf{Less code to write.} You describe \emph{what} you
        want, not \emph{how} to do it. \texttt{map[key] = value}
        hides a hash computation, bucket lookup, collision handling,
        and possible rehash.
  \item \textbf{Less code to read.} A colleague scanning your code
        sees the ADT's operations and immediately knows the shape of
        the data flow -- ``oh, it's a queue, so FIFO'' -- without
        studying a custom linked-list implementation.
  \item \textbf{Less code to debug.} Bugs in the ADT's implementation
        live in one place; bugs in your algorithm live in yours. The
        two don't mix.
\end{itemize}

\begin{examplebox}[Same job, two abstraction levels]
\begin{lstlisting}[basicstyle=\ttfamily\small,frame=single,language=C++]
// Low-level: roll our own LIFO using a raw array + top index.
int stack[100];
int top = -1;
stack[++top] = 3;
stack[++top] = 4;
int x = stack[top--];          // 4
int y = stack[top--];          // 3

// High-level: lean on the stack ADT from the STL.
#include <stack>
std::stack<int> s;
s.push(3);
s.push(4);
int x = s.top(); s.pop();       // 4
int y = s.top(); s.pop();       // 3
\end{lstlisting}
Both snippets do the same thing. The first mixes the algorithm with
the storage. The second talks only in stack operations -- push,
top, pop -- so a reader grasps the intent in one pass.
\end{examplebox}

\begin{warnbox}[Abstraction doesn't erase cost]
The textbook is explicit about this: \emph{``knowledge of the
underlying implementation is needed to analyze or improve the
runtime efficiency.''} A \texttt{std::map} and a
\texttt{std::unordered\_map} expose the same dictionary ADT, but one
is $O(\log n)$ and the other is expected $O(1)$ per operation --
with different worst-case behavior, cache profiles, and memory
overhead. When performance matters, the structure matters. Pretending
otherwise is how teams end up with ``mysteriously slow'' code.
\end{warnbox}

\subsection{ADTs in standard libraries}

Almost every modern language ships an \emph{implementation} of the
common ADTs as part of its standard library. You rarely implement a
queue or a hash map from scratch in production; you use the
language's provided one.

\begin{center}
\footnotesize
\renewcommand{\arraystretch}{1.3}
\begin{tabular}{@{}p{2.2cm}p{3.6cm}p{7.1cm}@{}}
\hline
\textbf{Language} & \textbf{Library} & \textbf{ADTs you get} \\
\hline
C++    & Standard Template Library (STL) &
        \texttt{vector}, \texttt{list}, \texttt{deque},
        \texttt{stack}, \texttt{queue}, \texttt{priority\_queue},
        \texttt{set}, \texttt{map}, \texttt{unordered\_set},
        \texttt{unordered\_map} \\
Python & Built-ins + \texttt{collections} &
        \texttt{list}, \texttt{tuple}, \texttt{set}, \texttt{dict},
        \texttt{deque}, \texttt{defaultdict}, \texttt{Counter},
        \texttt{heapq} \\
Java   & Java Collections Framework (JCF) &
        \texttt{ArrayList}, \texttt{LinkedList}, \texttt{HashSet},
        \texttt{TreeSet}, \texttt{HashMap}, \texttt{TreeMap},
        \texttt{ArrayDeque}, \texttt{PriorityQueue} \\
\hline
\end{tabular}
\end{center}

\begin{notebox}[How much do languages let you choose?]
Languages differ in how much control they give you over the
underlying structure:
\begin{itemize}[leftmargin=*]
  \item \textbf{C++} is the most explicit. Adaptors like
        \texttt{std::stack} and \texttt{std::priority\_queue} accept
        the underlying container as a template parameter; separate
        types (\texttt{map} vs \texttt{unordered\_map},
        \texttt{list} vs \texttt{forward\_list}) make the choice a
        deliberate decision.
  \item \textbf{Python}'s \texttt{list} is always a dynamic array;
        \texttt{dict} is always a hash table. The runtime picks; you
        don't.
  \item \textbf{Java} exposes interface-vs-implementation explicitly:
        you declare a variable as \texttt{List<String>} and assign
        it either an \texttt{ArrayList} or a \texttt{LinkedList}.
\end{itemize}
\end{notebox}

\subsection{Why CS 300 still has you build them by hand}

Production code uses the standard library. This course has you
implement linked lists, hash tables, heaps, and trees from scratch.
That's deliberate, and for two reasons.

\begin{ideabox}[Building vs.\ using]
\begin{enumerate}[leftmargin=*]
  \item \textbf{Understanding.} To reason about performance you have
        to know what is happening under the interface. Implementing
        a hash table once makes expressions like ``good hash
        function,'' ``load factor,'' and ``worst-case $O(n)$ on
        degenerate input'' stop being vocabulary and start being
        intuition.
  \item \textbf{Design judgment.} Choosing between
        \texttt{std::map} and \texttt{std::unordered\_map} in a real
        project requires knowing what each has to do internally.
        Library knowledge without structural knowledge leads to
        mysterious performance cliffs.
\end{enumerate}
\end{ideabox}

\begin{warnbox}[In real projects, reach for the library first]
Outside this course, re-implementing a hash map is a bright red flag
in a code review. The standard library's versions are tuned over
decades, heavily tested, and portable. Build-your-own is for
\emph{learning the idea}, not \emph{shipping the product}.
\end{warnbox}

\subsection{Choosing the right ADT is the actual skill}

Because the common ADTs are available off the shelf, the hard part
isn't implementing them -- it's recognizing which one the problem
in front of you is asking for.

\begin{center}
\footnotesize
\renewcommand{\arraystretch}{1.25}
\begin{tabular}{@{}p{6.2cm}p{5.2cm}@{}}
\hline
\textbf{Problem shape} & \textbf{ADT} \\
\hline
``I need the most recent item first.''          & Stack \\
``I need the oldest waiting item first.''       & Queue \\
``I need the \emph{most important} item first.'' & Priority queue \\
``I need fast lookup by a key.''                 & Map / dictionary \\
``I need a collection of distinct items, order doesn't matter.'' & Set \\
``I need ordered access by index.''              & Dynamic array / list \\
\hline
\end{tabular}
\end{center}

\begin{notebox}[Where this is going]
This closes the ``setup'' arc of the course. Chapter 4 starts the
walk: pick an ADT or structure, build it, analyze it, use it. The
background from 3.1--3.5 -- bits, the seven structures, the
algorithm / structure relationship, ADTs, and standard libraries --
is the common language every later chapter speaks.
\end{notebox}

\section{3.6 Searching and Algorithm Runtime}

This section restarts the algorithms thread now that we have a
vocabulary of data structures to search through. The worked example
is \emph{linear search}, the simplest search algorithm there is. The
real payoff of the section, though, is the first careful look at
\textbf{runtime as a function of input size} -- the idea that turns
into Big-O in \textsection 3.10.

\subsection{Linear search}

\begin{defnbox}[Linear search]
\textbf{Linear search} (also called sequential search) walks a
collection from the first element to the last, comparing each
element to a search key, returning the index where it matches or a
sentinel value (often \texttt{-1}) if it never does.
\end{defnbox}

\begin{examplebox}[Linear search in C++]
\begin{lstlisting}[basicstyle=\ttfamily\small,frame=single,language=C++]
int linearSearch(const std::vector<int>& nums, int key) {
    for (std::size_t i = 0; i < nums.size(); ++i) {
        if (nums[i] == key) {
            return static_cast<int>(i);
        }
    }
    return -1;  // not found
}

int main() {
    std::vector<int> nums = {2, 4, 7, 10, 11, 32, 45, 87};
    int idx = linearSearch(nums, 10);
    if (idx == -1) std::cout << "not found\n";
    else           std::cout << "found at index " << idx << '\n';
}
\end{lstlisting}
\end{examplebox}

\begin{notebox}[STL already has it]
\texttt{std::find(v.begin(), v.end(), key)} returns an iterator to
the first matching element, or \texttt{v.end()} if none. Same
algorithm, idiomatic C++. Hand-rolling the loop is a learning
exercise; in production, use the library.
\end{notebox}

\subsection{Why linear -- and when it's fine}

Linear search makes \emph{no} assumptions about the data. It works
on an unsorted array, a linked list, or any sequence you can iterate
once. That's the whole point: it's the baseline.

\begin{itemize}[leftmargin=*]
  \item \textbf{Best case.} Key is at index 0 -- \textbf{1 comparison}.
  \item \textbf{Worst case.} Key is last, or absent -- \textbf{$n$ comparisons}.
  \item \textbf{Average case.} On uniformly random hit positions,
        roughly $n/2$ comparisons. When the key is absent, always $n$.
\end{itemize}

\begin{warnbox}[Linear search on a linked list is still $O(n)$]
Don't confuse \emph{structure} with \emph{algorithm}. Linear search
is $O(n)$ whether the container is a \texttt{std::vector} or a
\texttt{std::list}. The linked list's poor random access doesn't
make linear search any worse -- linear search only ever moves to the
\emph{next} element, which is cheap for both structures.
\end{warnbox}

\subsection{Runtime as a function of input size}

The key idea the textbook introduces here: runtime grows with the
size of the input. Counting \emph{actual seconds} depends on the
machine. Counting \emph{the number of operations the algorithm does}
is a machine-independent proxy that predicts the seconds.

\begin{ideabox}[Order of growth, informally]
If an algorithm does ``about $N$'' operations on an input of size
$N$, we say its runtime is \textbf{on the order of} $N$. Doubling
the input doubles the work. Linear search is on the order of $N$;
we'll soon call this $O(N)$.
\end{ideabox}

\subsection*{Concrete numbers}

The scale at which ``$O(N)$'' starts to feel slow surprises people
the first time they see it. Suppose one comparison takes 1
microsecond ($10^{-6}$ s).

\begin{center}
\footnotesize
\renewcommand{\arraystretch}{1.25}
\begin{tabular}{@{}rll@{}}
\hline
\textbf{$N$} & \textbf{Worst-case comparisons} & \textbf{Wall time @ 1\,\textmu s each} \\
\hline
$10^3$               & 1{,}000        & 0.001 s \\
$10^6$               & 1{,}000{,}000  & 1 s \\
$10^7$               & $10^7$         & 10 s \\
$2 \times 10^8$ (Amazon's catalog, ish) & $2 \times 10^8$ & 200 s $\approx$ 3 min 20 s \\
$10^9$               & $10^9$         & about 17 min \\
\hline
\end{tabular}
\end{center}

\begin{warnbox}[``Fast enough'' is always ``fast enough \emph{for what?}'']
1 second is fine if you're searching once. For an interactive search
box that fires on every keystroke, it's catastrophic. For a
background indexer, it may be fine. \emph{Runtime is a constraint,
not a verdict -- what matters is whether it fits the use case.}
\end{warnbox}

\subsection{Ignoring constants: why we say ``on the order of $N$''}

If your linear search does 3 operations per element (bounds check,
array access, comparison) instead of 1, the constant changes but the
\emph{shape} of the growth does not: doubling $N$ still doubles the
work.

\begin{ideabox}[The shape is what matters]
For large $N$, the \emph{shape} of the growth (constant, linear,
quadratic, logarithmic \ldots) dominates any constant-factor
difference. An $O(N)$ algorithm with a big constant will eventually
beat an $O(N^2)$ algorithm with a small one. ``Eventually'' may be
sooner than you think: at $N = 10{,}000$, a $100N$ algorithm does
$10^6$ ops; a $0.1 N^2$ algorithm does $10^7$. The shape has already
won.
\end{ideabox}

\subsection{Searching, structures, and algorithms together}

The textbook frames this section as ``algorithms,'' but notice how
much depends on the structure underneath. Linear search is the
universal baseline because it makes no assumptions. If you know more
about the structure, you can do much better:

\begin{center}
\footnotesize
\renewcommand{\arraystretch}{1.25}
\begin{tabular}{@{}p{5.5cm}p{4.5cm}p{2.8cm}@{}}
\hline
\textbf{What you know about the data} & \textbf{Algorithm} & \textbf{Worst case} \\
\hline
Nothing; unsorted sequence                & Linear search      & $O(N)$ \\
Array is sorted                           & Binary search      & $O(\log N)$ \\
Data lives in a hash table keyed by what you're searching for & Hash lookup & $O(1)$ expected \\
Data lives in a balanced BST keyed by the search value & BST lookup & $O(\log N)$ \\
\hline
\end{tabular}
\end{center}

\begin{notebox}[Where this is going]
\textsection 3.10 formalizes the ``on the order of'' idea into Big-O
notation, and \textsection 3.8 introduces binary search -- the first
algorithm that does substantially better than $O(N)$, but only on a
\emph{sorted} structure. That's the cycle you'll see again and again:
pick an algorithm, pay the structure's cost; pick a structure, unlock
a better algorithm.
\end{notebox}

\section{3.8 Binary Search}

Section 3.6 covered linear search -- the universal baseline. This
section introduces the first algorithm that does substantially better:
\textbf{binary search}. The catch, and the whole point, is that it
only works if the data is \emph{sorted and randomly accessible}. That
trade-off -- accept a precondition on the structure, unlock a much
faster algorithm -- is a pattern you will see repeatedly for the
rest of the course.

\subsection{The idea}

Think about how you find a name in a physical phone book. You don't
read from the beginning. You open near the middle, see whether your
target comes before or after, and immediately ignore the half that
doesn't contain it. Repeat.

\begin{defnbox}[Binary search]
\textbf{Binary search} finds an item in a \emph{sorted, randomly
accessible} sequence by repeatedly looking at the middle element and
discarding the half that cannot contain the key. Each step halves the
search space; the algorithm terminates when the key is found or the
space is empty.
\end{defnbox}

\begin{warnbox}[Two preconditions -- both matter]
Binary search requires both:
\begin{itemize}[leftmargin=*]
  \item \textbf{Sorted.} Without order, ``look at the middle'' tells
        you nothing about which half the key is in.
  \item \textbf{Random access in \(O(1)\).} You have to jump straight
        to the middle. On an array, that's \texttt{a[mid]}. On a
        \emph{linked list}, reaching the middle takes $O(n)$ -- and
        that extra cost destroys the $O(\log n)$ advantage.
\end{itemize}
So binary search on a \texttt{std::list} is \emph{not} a thing; on
\texttt{std::vector} (or a plain array), it is.
\end{warnbox}

\subsection{Binary search in C++}

\begin{examplebox}[Iterative binary search]
\begin{lstlisting}[basicstyle=\ttfamily\small,frame=single,language=C++]
int binarySearch(const std::vector<int>& nums, int key) {
    int low = 0;
    int high = static_cast<int>(nums.size()) - 1;

    while (low <= high) {
        int mid = low + (high - low) / 2;   // overflow-safe midpoint
        if      (nums[mid] == key) return mid;
        else if (nums[mid]  < key) low  = mid + 1;
        else                       high = mid - 1;
    }
    return -1;  // not found
}
\end{lstlisting}
The textbook writes \texttt{(low + high) / 2}. That formula can
overflow when \texttt{low + high} exceeds \texttt{INT\_MAX}. The
idiomatic fix is \texttt{low + (high - low) / 2}, which computes the
same midpoint without the intermediate sum blowing up.
\end{examplebox}

\begin{notebox}[Real-world gotcha, not a hypothetical]
This is a famous bug. JDK's \texttt{Arrays.binarySearch} had the
na\"ive formula for nearly a decade before Josh Bloch documented it
(``Extra, Extra -- Read All About It: Nearly All Binary Searches and
Mergesorts are Broken,'' 2006). Use the overflow-safe form.
\end{notebox}

\subsection*{Recursive binary search}

Binary search is also the textbook example of a recursive algorithm
that calls itself on a \emph{shrinking} problem.

\begin{examplebox}[Recursive binary search]
\begin{lstlisting}[basicstyle=\ttfamily\small,frame=single,language=C++]
int binarySearch(const std::vector<int>& nums,
                 int key, int low, int high) {
    if (low > high) return -1;                 // base case
    int mid = low + (high - low) / 2;
    if (nums[mid] == key) return mid;
    if (nums[mid]  < key) return binarySearch(nums, key, mid + 1, high);
    return                   binarySearch(nums, key, low, mid - 1);
}
\end{lstlisting}
Same algorithm; two views. The iterative version updates a window;
the recursive version calls itself on a smaller window. The iterative
version is usually preferred in C++ because it avoids stack frames.
\end{examplebox}

\subsection{Why $O(\log n)$}

\begin{ideabox}[The halving argument]
Start with $n$ candidates. After one step, at most $n/2$. After two
steps, $n/4$. After $k$ steps, $n / 2^k$. The loop ends when the
window is empty; that happens when $n / 2^k \le 1$, i.e.\ when
$k \ge \log_2 n$.

Binary search does at most
\[
    \lfloor \log_2 n \rfloor + 1
\]
comparisons in the worst case. The runtime is \(\boldsymbol{O(\log n)}\).
\end{ideabox}

\subsection*{Why $\log n$ is so dramatic}

The concrete comparison with linear search is what sells binary
search. One comparison at 1\,\textmu s:

\begin{center}
\footnotesize
\renewcommand{\arraystretch}{1.25}
\begin{tabular}{@{}rrrr@{}}
\hline
\textbf{$N$} & \textbf{Linear worst} & \textbf{Binary worst} & \textbf{Ratio} \\
\hline
$10^3$          & 1{,}000 ops  (1 ms)      & 10 ops (10\,\textmu s)       & 100$\times$ \\
$10^6$          & $10^6$ ops (1 s)         & 20 ops (20\,\textmu s)       & 50{,}000$\times$ \\
$2 \times 10^8$ & $2 \times 10^8$ (3+ min) & 28 ops (28\,\textmu s)       & $\approx 7 \times 10^6$$\times$ \\
$10^9$          & $10^9$ ops ($\approx$17 min) & 30 ops (30\,\textmu s)   & $\approx 3 \times 10^7$$\times$ \\
\hline
\end{tabular}
\end{center}

\begin{ideabox}[The ``log bar''  for intuition]
$\log_2 n$ is the number of times you can halve $n$ before reaching
$1$. For any $n$ up to a billion, that's at most 30. Double the
input, you add \emph{one} operation. This is the first example in the
course of truly different growth behavior -- the gap between linear
and logarithmic is not ``a bit faster,'' it's categorical.
\end{ideabox}

\subsection{The bigger lesson}

\begin{warnbox}[Sorting isn't free]
If your data isn't already sorted, you have to sort it first --
which is $O(n \log n)$ for the first sort. One binary search on a
freshly sorted array does not beat one linear search. Binary search
pays off when you will do \emph{many} searches against the same
sorted data, amortizing the sort cost across the searches.
\end{warnbox}

\begin{ideabox}[Structure + precondition = algorithm]
This is the first of many ``precondition unlocks algorithm''
patterns you will meet. The full arc:
\begin{itemize}[leftmargin=*]
  \item \emph{Unsorted sequence.} Linear search, $O(n)$.
  \item \emph{Sorted sequence, random access.} Binary search, $O(\log n)$.
  \item \emph{Binary search tree.} Tree descent, $O(\log n)$ expected.
  \item \emph{Hash table keyed by the search value.} Direct lookup, $O(1)$ expected.
\end{itemize}
Each improvement comes from structuring the data to answer the
question you'll ask. Algorithm design is in large part the art of
preparing the data so that lookup is cheap.
\end{ideabox}

\begin{notebox}[STL ships binary search]
\texttt{std::binary\_search(v.begin(), v.end(), key)} returns a
\texttt{bool}. \texttt{std::lower\_bound} returns the iterator to the
first element $\ge$ key (where you'd insert). You'll use these a lot
when a sorted structure is the right choice.
\end{notebox}

\section{3.9 Constant-Time Operations}

Binary search was the first algorithm where the \emph{number of
operations} mattered more than the speed of the processor. This
section makes precise what ``one operation'' even means for runtime
analysis. The short answer: we count \textbf{constant-time
operations} -- the unit of work the CPU can do in a fixed time
independent of input values.

\begin{defnbox}[Constant-time operation]
A \textbf{constant-time operation} is one whose execution time does
not depend on the size of the input data. The same operation takes
roughly the same amount of time whether the operand is small or
large. We call this $O(1)$.
\end{defnbox}

\begin{ideabox}[Why we count operations instead of seconds]
Seconds change with the machine. Operations don't. If an algorithm
does $n$ constant-time operations, then on a 2$\times$ faster CPU it
takes half as long -- but it still does $n$ operations. Counting
operations gives us a measure of an algorithm that is \emph{about
the algorithm}, not about the hardware.
\end{ideabox}

\subsection{What counts as constant time}

The short list of things that are typically $O(1)$ on a modern
processor:

\begin{center}
\footnotesize
\renewcommand{\arraystretch}{1.3}
\begin{tabular}{@{}p{6cm}p{6.8cm}@{}}
\hline
\textbf{Operation} & \textbf{Why it's constant} \\
\hline
Arithmetic on fixed-size numeric types (\texttt{int}, \texttt{double}) & Hardware does it in a bounded number of cycles regardless of values. \\
Assignment of a fixed-size value (scalar, pointer, reference) & Copies a fixed number of bytes. \\
Comparison of fixed-size values                            & One CPU instruction. \\
Indexed access to an array element (\texttt{a[i]})         & One pointer arithmetic + one memory load. \\
Push / pop on a \texttt{std::vector} (amortized)           & Constant on average; occasional resize is paid for across many pushes. \\
Hash-table lookup (\emph{expected})                        & One hash computation + one bucket probe on average. \\
Dereferencing a pointer                                    & One memory load. \\
Reading the length stored by a container (\texttt{v.size()}) & Returns a cached value. \\
\hline
\end{tabular}
\end{center}

\begin{examplebox}[Operations that \emph{look} simple but aren't constant time]
\begin{lstlisting}[basicstyle=\ttfamily\small,frame=single,language=C++]
std::string a(1'000'000, 'x');
std::string b(1'000'000, 'y');

std::string c = a + b;          // O(n) -- allocates and copies
std::string d = a;              // O(n) -- copy constructor
bool same = (a == b);           // O(n) -- compares every character
int count = std::count(v.begin(), v.end(), 3);  // O(n)
\end{lstlisting}
The \texttt{+}, \texttt{=}, and \texttt{==} operators look like one
operation at the language level, but they walk every character of
every string. The source code has constant width; the runtime does
not.
\end{examplebox}

\begin{warnbox}[``It's one line'' $\neq$ ``it's $O(1)$'']
The biggest beginner mistake in complexity analysis is confusing
\emph{syntax} with \emph{semantics}. Any operation on a container
whose size grows with input -- concatenation, copying, containment
checks, sort, reverse, clear of a large structure -- has a cost
proportional to that size. Even when the call site is a single
token, the work isn't free.
\end{warnbox}

\subsection{Why fixed size matters}

Notice the repeating phrase: ``fixed-size integer or floating-point
values.'' A modern CPU's adder has a fixed number of bits -- 32 or
64 -- and its time-per-add doesn't depend on whether the operands
are 3 or 2{,}147{,}483{,}646.

\begin{ideabox}[Arbitrary-precision arithmetic is \emph{not} $O(1)$]
Python's built-in integers are arbitrary-precision: you can compute
$10^{10000}$ directly. The trade-off is that arithmetic on them is
no longer constant-time. Adding two $n$-digit numbers is $O(n)$;
multiplying them is $O(n \log n)$ with fast algorithms, $O(n^2)$ with
the schoolbook method.

C++ \texttt{int}, \texttt{long long}, and \texttt{double} are
fixed-size by design. Arithmetic on them is genuinely constant-time
-- at the cost of overflow being possible. Different languages, same
trade-off: you get either $O(1)$ arithmetic \emph{or} unlimited
precision, not both.
\end{ideabox}

\subsection{Counting the right things}

In runtime analysis we don't count every line of code. We pick a
representative operation and count how many times it runs.

\begin{examplebox}[What to count, by algorithm]
\begin{itemize}[leftmargin=*]
  \item \textbf{Linear search.} Count \emph{comparisons} between the
        search key and array elements. Each comparison is $O(1)$, so
        total work is proportional to the number of comparisons.
  \item \textbf{Sorting algorithms.} Count comparisons \emph{and/or}
        swaps; both are constant-time.
  \item \textbf{Binary search.} Count iterations (or recursive
        calls); each iteration does constant work.
\end{itemize}
\end{examplebox}

\subsection*{The hidden-loop trap}

\begin{warnbox}[Watch for hidden loops inside library calls]
\begin{lstlisting}[basicstyle=\ttfamily\small,frame=single,language=C++]
for (std::size_t i = 0; i < n; ++i) {
    std::string msg = "item " + std::to_string(i);  // O(len)
    output += msg;                                   // O(|output|)
}
\end{lstlisting}
This is \textbf{not} $O(n)$. The \texttt{output += msg} concatenation
grows more expensive each iteration because \texttt{|output|} keeps
growing. Total work is $O(n^2)$. A reader counting ``one loop, so
linear'' gets the answer spectacularly wrong.
\end{warnbox}

\subsection{The rule of thumb}

\begin{ideabox}[The ``one at a time'' heuristic]
If an operation touches a \emph{constant} amount of memory -- one
integer, one pointer, one cache line, one array cell -- it's almost
certainly $O(1)$. If it touches memory whose size grows with the
input, assume it isn't. When in doubt:
\begin{enumerate}[leftmargin=*]
  \item Does it allocate?
  \item Does it copy?
  \item Does it call another function that loops?
\end{enumerate}
Any ``yes'' is a hint that the operation is not constant-time.
\end{ideabox}

\begin{notebox}[Where this is going]
This section defined the unit of measurement. The next section
(3.10, ``Growth of functions and complexity'') takes the count of
constant-time operations -- as a function of input size -- and
classifies algorithms by how that function grows. That classification
is Big-O notation, formalized in the next section, \textsection 3.10.
\end{notebox}

\section{3.10 Growth of Functions and Complexity}

Section 3.9 defined the unit of measurement: constant-time operations.
This section asks \emph{how the count of those operations grows as
the input grows}. That growth is what we summarize with Big-O (and
its cousins $\Omega$, $\Theta$). These are the three core pieces of
vocabulary that will structure every runtime claim for the rest of
the course.

\subsection{Best case, worst case, and $T(N)$}

For an algorithm running on input of size $N$, the number of
constant-time operations it performs is a function of $N$. We'll
call that function $T(N)$. Usually $T(N)$ depends on more than just
$N$ -- it depends on which particular input of size $N$ we're running
on. Hence the standard two-sided view:

\begin{defnbox}[Best / worst case runtime]
\textbf{Best-case $T(N)$} is the fewest constant-time operations the
algorithm could do on any input of size $N$. \textbf{Worst-case
$T(N)$} is the most. Average-case is an in-between notion that
requires a probability distribution over inputs.
\end{defnbox}

\begin{examplebox}[Linear search as a two-sided picture]
\begin{itemize}[leftmargin=*]
  \item Best case: key is at index 0, so $T_{\text{best}}(N) = 1$.
  \item Worst case: key is last, or absent, so $T_{\text{worst}}(N) = N$.
\end{itemize}
Both are true statements about the same algorithm. You report
whichever matches the guarantee you're after.
\end{examplebox}

\subsection{Upper and lower bounds}

Rather than carrying the exact $T(N)$ formula around (with all its
constants and lower-order terms), we bound it with simpler functions.

\begin{defnbox}[Upper / lower bound]
A function $f(N)$ is an \textbf{upper bound} on $T(N)$ if, for all
sufficiently large $N$, $T(N) \le f(N)$. It's a \textbf{lower bound}
if $T(N) \ge f(N)$ for all sufficiently large $N$.
\end{defnbox}

\begin{examplebox}[Bounds for a messy $T(N)$]
Suppose an algorithm's worst-case runtime is
\[
  T(N) = 3N^2 + 10N + 17.
\]
\begin{itemize}[leftmargin=*]
  \item $f(N) = 30 N^2$ is an upper bound. ($30 N^2 \ge 3N^2 + 10N + 17$
        for all $N \ge 1$.)
  \item $f(N) = N^2$ is a lower bound.
\end{itemize}
Either bound gives a cleaner statement than the exact $T(N)$, at the
cost of some precision.
\end{examplebox}

\subsection{Asymptotic notation: strip the constants}

For cross-algorithm comparison we don't care about the constants in
front of $N^2$. We care that it's an $N^2$. That's what Big-O,
$\Omega$, and $\Theta$ formalize: \emph{rate of growth, constants and
lower-order terms absorbed}.

\begin{defnbox}[Asymptotic notations]
Let $T(N)$ be an algorithm's runtime (in constant-time operations)
and $f(N)$ a simpler reference function. Then:
\begin{itemize}[leftmargin=*]
  \item $T(N) = O(f(N))$ means there's a positive constant $c$ with
        $T(N) \le c \cdot f(N)$ for all $N \ge N_0$. \\
        \emph{Big-O is an upper bound on growth rate.}
  \item $T(N) = \Omega(f(N))$ means there's a positive constant $c$
        with $T(N) \ge c \cdot f(N)$ for all $N \ge N_0$. \\
        \emph{Big-$\Omega$ is a lower bound on growth rate.}
  \item $T(N) = \Theta(f(N))$ means $T(N)$ is both $O(f(N))$ and
        $\Omega(f(N))$ -- the two bounds meet. \\
        \emph{Big-$\Theta$ is a tight bound.}
\end{itemize}
\end{defnbox}

\begin{ideabox}[The mental picture]
\begin{itemize}[leftmargin=*]
  \item \textbf{$O$}: ``The runtime doesn't grow faster than this.''
  \item \textbf{$\Omega$}: ``The runtime doesn't grow slower than this.''
  \item \textbf{$\Theta$}: ``The runtime grows exactly like this.''
\end{itemize}
$O$ is what you say when you want to promise \emph{it won't blow up};
$\Omega$ is what you say to argue \emph{it can't be better than this};
$\Theta$ is what you say when you've pinned it from both sides.
\end{ideabox}

\begin{notebox}[Little-oh and little-omega: when the bound is strictly looser]
$O$ and $\Omega$ may or may not be \emph{tight} -- $2n = O(n^2)$ is
true (linear grows no faster than quadratic), but it's also a loose
bound. CLRS introduces two more notations for the strictly-looser
case:

\begin{itemize}[leftmargin=*,noitemsep]
  \item $f(n) = o(g(n))$ (\textbf{little-oh}): for \emph{every}
        positive constant $c$, there exists $n_0$ such that
        $f(n) < c \cdot g(n)$ for all $n \ge n_0$. Equivalently,
        $\lim_{n\to\infty} f(n)/g(n) = 0$. Strict upper bound.
  \item $f(n) = \omega(g(n))$ (\textbf{little-omega}): the mirror
        image. $\lim_{n\to\infty} f(n)/g(n) = \infty$. Strict lower bound.
\end{itemize}

The cleanest mental model is the analogy with real-number comparisons:

\begin{center}
\small
\begin{tabular}{@{}ll@{}}
$f = O(g)$       & like $f \le g$ \\
$f = \Omega(g)$  & like $f \ge g$ \\
$f = \Theta(g)$  & like $f = g$ \\
$f = o(g)$       & like $f < g$ \quad (strict) \\
$f = \omega(g)$  & like $f > g$ \quad (strict) \\
\end{tabular}
\normalsize
\end{center}

So $2n = O(n^2)$ AND $2n = o(n^2)$ -- the latter is the precise
statement (linear is \emph{strictly} smaller than quadratic).
$2n = O(n)$ is also true, but $2n \ne o(n)$ -- linear is not
strictly smaller than itself. Use $o$ and $\omega$ when you want to
signal ``definitely not just a loose bound.''
\end{notebox}

\subsection*{What $c$ and $N_0$ buy you}

The constant $c$ absorbs the hardware, the language, and any
constant-factor overhead. The threshold $N_0$ lets us \emph{ignore
small inputs}, where lower-order terms and one-time setup costs
distort the picture.

\begin{examplebox}[Why the constants go away]
Let $T(N) = 3N^2 + 10N + 17$. Then $T(N) = O(N^2)$ because for
$c = 30$ and $N \ge 1$,
\[
  3N^2 + 10N + 17 \;\le\; 3N^2 + 10N^2 + 17N^2 \;=\; 30N^2.
\]
We gave up the exact coefficient ($3$) to get a clean, hardware-free
statement: this algorithm's runtime \emph{grows like $N^2$}.
\end{examplebox}

\subsection{How to read a big-O claim}

\begin{warnbox}[Common traps]
\begin{itemize}[leftmargin=*]
  \item \textbf{Big-O is \emph{not} the ``worst case.''} You can have
        Big-O of \emph{any} case -- best, worst, average. ``Linear
        search worst case is $O(N)$'' is a bound on the worst case.
        ``Linear search best case is $O(1)$'' is a bound on the best
        case. Both are legitimate.
  \item \textbf{Big-O doesn't promise tightness.} Saying ``linear
        search worst case is $O(N^2)$'' is technically true -- $N$
        grows no faster than $N^2$. It's just a loose upper bound.
        That's why we prefer $\Theta$ when we can prove it.
  \item \textbf{Constants and lower-order terms are not mysteries.}
        They're real; $100N$ is slower than $N$ in practice. Big-O is
        a coarse tool for comparing \emph{shapes}, not a license to
        ignore the constant when performance matters.
\end{itemize}
\end{warnbox}

\subsection{The growth-rate hierarchy}

For large $N$, the following functions are ordered by how fast they
grow. Every entry dominates the ones above it.

\begin{center}
\footnotesize
\renewcommand{\arraystretch}{1.3}
\begin{tabular}{@{}lll@{}}
\hline
\textbf{Class} & \textbf{Name}             & \textbf{Examples} \\
\hline
$O(1)$         & Constant                  & Array index, hash lookup \\
$O(\log N)$    & Logarithmic               & Binary search, balanced-BST lookup \\
$O(N)$         & Linear                    & Linear search, one pass over an array \\
$O(N \log N)$  & Linearithmic / ``$N$ log $N$'' & Merge sort, heap sort, good sorts \\
$O(N^2)$       & Quadratic                 & Bubble sort, insertion sort worst case, nested loops \\
$O(N^3)$       & Cubic                     & Naive matrix multiply, certain DP algorithms \\
$O(2^N)$       & Exponential               & Brute-force subset enumeration \\
$O(N!)$        & Factorial                 & Brute-force permutations (traveling salesman) \\
\hline
\end{tabular}
\end{center}

\begin{ideabox}[The gaps between classes are enormous]
\footnotesize
\renewcommand{\arraystretch}{1.2}
At $N = 100$ (only!), the difference between the classes is already
wild:
\[
\begin{array}{r|r}
\text{Class} & \text{ops at } N=100 \\ \hline
O(1)           & 1 \\
O(\log N)      & 7 \\
O(N)           & 100 \\
O(N \log N)    & \approx 664 \\
O(N^2)         & 10{,}000 \\
O(2^N)         & \approx 1.27 \times 10^{30} \\
O(N!)          & \approx 9.33 \times 10^{157} \\
\end{array}
\]
The jump from quadratic to exponential is where problems go from
``slow'' to ``literally cannot be solved in the lifetime of the
universe for modest $N$.''
\end{ideabox}

\subsection{In practice}

\begin{notebox}[What you'll actually say in code reviews]
In day-to-day work you almost always state runtime as Big-O of the
worst case: ``\texttt{std::unordered\_map::find} is $O(1)$ expected,''
``your nested loop is $O(N^2)$,'' ``this sort is $O(N \log N)$.''
$\Omega$ and $\Theta$ show up in analysis problems and lower-bound
arguments more than in casual engineering chatter -- but knowing
them lets you read correct statements like ``no comparison sort can
do better than $\Omega(N \log N)$.''
\end{notebox}

\begin{ideabox}[Halfway recap -- before we dive into sorts]
You've now been through the first half of ch.~3. Before sorting algorithms
start piling up, consolidate:

\medskip
\textbf{What you've built (3.1-3.10):}
\begin{enumerate}
    \item \textbf{Bits \& ASCII.} Data in memory is bits; characters are
    just small integers with an agreed table. You can move between
    \texttt{char}, \texttt{int}, and \texttt{cctype} classification
    without friction.
    \item \textbf{Seven data structures named.} Array, linked list, tree,
    BST, heap, hash table, graph. You can't implement them all yet, but
    you know their names and what shape of problem each is good for.
    \item \textbf{ADT vs.\ data structure.} An ADT is the \emph{contract}
    (``a stack has push/pop/top'' — no mention of \emph{how}). A data
    structure is the \emph{implementation}. Every C++ STL container is
    an ADT realized by one or more data structures
    (\texttt{std::stack<T>} is an ADT; its default implementation is
    \texttt{std::deque<T>}).
    \item \textbf{Two searches contrasted.} Linear search is $O(n)$ and
    works on any data. Binary search is $O(\log n)$ but requires sorted
    data — which motivates the next half of the chapter.
    \item \textbf{Big-O as a working tool.} You can read ``constant
    time,'' ``logarithmic,'' ``linear,'' ``quadratic,'' ``exponential''
    and point at the growth-rate hierarchy.
\end{enumerate}

\medskip
\textbf{The through-line into sorting.} Binary search is the payoff of
sorted data. That makes ``how expensive is sorting?'' the immediate
follow-up question: if sorting costs $O(n \log n)$ and then you do one
search at $O(\log n)$, you've paid more than linear search ($O(n)$) for
one lookup — so sorting only pays when you search \emph{many} times.
Every sort in the next half of the chapter is a different attempt to make
that one-time cost as small as possible.

\medskip
\textbf{What's ahead (3.14-3.21):} seven sorting algorithms plus a fast-sort
overview. They split into three families:
\begin{itemize}
    \item \textbf{Incremental:} bubble (covered in ch\_13 extras as pedagogy-only), insertion, selection, shell —
    easy to code, $O(n^2)$ (except shell).
    \item \textbf{Divide-and-conquer:} merge sort, quicksort — $O(n \log n)$
    average; the ones you'll actually ship.
    \item \textbf{Non-comparison:} radix sort — $O(n k)$, fast when $k$
    (key width) is small.
\end{itemize}

\medskip
\textbf{If any of 3.1-3.10 still feels shaky, stop here and re-read.} The
sorting analyses that follow all lean on the Big-O definition and the
growth-rate hierarchy. You don't need mastery yet; you need enough
fluency that ``this loop is $O(n)$'' and ``this outer-over-inner loop
is $O(n^2)$'' don't slow you down.
\end{ideabox}

\section{3.14 Sorting: Introduction}

Searching (3.6, 3.8) was the first algorithmic problem we analyzed.
\textbf{Sorting} is the second, and it's arguably the more important
one: binary search, priority queues, deduplication, many-pass merges,
and much of computational geometry all assume data you can sort
first. This short section sets the stage; the rest of this chapter
walks through insertion, selection, shell, quicksort, merge, and radix
sort. (Bubble sort lives in ch\_13's extras chapter, since it is
pedagogy-only; heap sort waits for ch\_7, once we have heaps.)

\begin{defnbox}[Sorting]
\textbf{Sorting} is the process of rearranging a sequence of elements
into non-decreasing (``ascending'') or non-increasing (``descending'')
order under some total ordering. For numbers, the order is $\le$;
for strings, usually lexicographic; for custom types, whatever
comparison function you supply.
\end{defnbox}

\subsection{Why sorting is its own field}

Sorting is the most-studied algorithmic problem in computer science.
It looks trivial at small sizes -- you could do it by hand -- and
turns out to be surprisingly deep.

\begin{ideabox}[What sorting unlocks]
Sorted data is cheaper to work with in several important ways.
\begin{itemize}[leftmargin=*]
  \item \textbf{Binary search.} $O(\log n)$ lookup instead of $O(n)$.
  \item \textbf{Duplicates.} Duplicates end up adjacent -- a single
        pass detects or removes them.
  \item \textbf{Set operations.} Intersection, union, and difference
        of two sorted sequences are each $O(n + m)$ (merge-style).
  \item \textbf{$k$-th smallest / largest.} Once sorted, it's an
        index lookup.
  \item \textbf{Human consumption.} Sorted output is the default
        contract of almost every report, table, and leaderboard.
\end{itemize}
Much of the time, the cost of sorting ($O(n \log n)$) is amortized
across many subsequent queries that are each much cheaper because
of the sort.
\end{ideabox}

\subsection{The fundamental constraint}

Why isn't sorting just ``put each element in its right slot''? Because
the program \emph{doesn't know} where the right slot is without
doing work to find out.

\begin{defnbox}[The access constraint]
A sorting algorithm can't see the whole list at once. It can only
\emph{observe or swap a small number of elements at a time} --
typically two, via a comparison and a conditional swap. Sorting
algorithms are the art of doing this efficiently.
\end{defnbox}

\begin{examplebox}[Swap as the atomic move]
\begin{lstlisting}[basicstyle=\ttfamily\small,frame=single,language=C++]
void swap(int& a, int& b) {
    int tmp = a;
    a = b;
    b = tmp;
}

// Or just:
std::swap(a, b);
\end{lstlisting}
Almost every comparison-based sort boils down to compare-two,
swap-if-needed, repeated in different patterns. Bubble sort scans
adjacent pairs. Selection sort finds the minimum of a suffix and
swaps it into place. Insertion sort shifts a value leftward through
an already-sorted prefix. Quicksort partitions around a pivot.
Merge sort combines sorted halves.
\end{examplebox}

\begin{notebox}[Why we can't ``just place each element in its slot'']
Placing each element in its final slot would mean \emph{already
knowing} where it goes, which means already having done the sort.
An algorithm starts knowing nothing beyond the ability to compare
two elements and swap two slots. The creativity in sorting algorithms
is in deciding which comparisons to make and in what order.
\end{notebox}

\subsection{What makes sorts differ}

When we compare two sorting algorithms we ask five questions:

\begin{center}
\footnotesize
\renewcommand{\arraystretch}{1.3}
\begin{tabular}{@{}p{3.2cm}p{9.2cm}@{}}
\hline
\textbf{Property} & \textbf{What it means} \\
\hline
Time complexity & Best / average / worst-case count of comparisons
                  and swaps as a function of $n$. \\
Space complexity & Extra memory beyond the input. \emph{In-place}
                  sorts use $O(1)$ extra; merge sort needs $O(n)$. \\
Stability       & Does the sort preserve the original order of
                  \emph{equal} elements? Useful when you've already
                  sorted by a secondary key. \\
Adaptivity      & Does it run faster on already-mostly-sorted input? \\
Access pattern  & Does it need random access (quicksort, binary
                  search during sort) or sequential only (merge sort)? \\
\hline
\end{tabular}
\end{center}

\begin{ideabox}[The big-picture result]
There's a proven lower bound: \emph{any comparison-based sort does
$\Omega(n \log n)$ comparisons in the worst case}. Merge sort, heap
sort, and well-implemented quicksort all meet this bound
asymptotically. Bubble sort, selection sort, and insertion sort do
not -- they're $\Theta(n^2)$. That gap is why ``toy sorts'' are
pedagogy and ``real sorts'' are engineering.
\end{ideabox}

\subsection*{Faster-than-comparison sorts}

\begin{notebox}[Specialized sorts can beat $n \log n$]
If you have more information about the values (say, they're small
non-negative integers, or short fixed-length keys), \emph{non-comparison}
sorts -- counting sort, radix sort, bucket sort -- can run in
$O(n)$ or $O(n + k)$. The $\Omega(n \log n)$ lower bound only applies
when the algorithm's only tool is comparison.
\end{notebox}

\subsection{In C++}

\begin{examplebox}[The library already sorts]
\begin{lstlisting}[basicstyle=\ttfamily\small,frame=single,language=C++]
#include <algorithm>
std::vector<int> v = {17, 3, 44, 6, 9};
std::sort(v.begin(), v.end());                     // ascending: 3 6 9 17 44
std::sort(v.begin(), v.end(), std::greater<int>()); // descending

std::stable_sort(v.begin(), v.end());               // preserves equal-order
std::partial_sort(v.begin(), v.begin() + 3, v.end()); // only first 3 sorted
\end{lstlisting}
\texttt{std::sort} is typically an \emph{introsort}: quicksort that
switches to heap sort when recursion gets too deep. Worst-case
$O(n \log n)$, average $O(n \log n)$ with a tiny constant. For
production code, start here.
\end{examplebox}

\begin{ideabox}[Where this is going]
The rest of the sorting material walks through the classic
algorithms: bubble, selection, and insertion sort (simple, quadratic,
pedagogical); then quicksort, merge sort, and heap sort ($n \log n$,
practical). Each is a different answer to the same question: ``given
only compare and swap, in what order should I issue them?''
Understanding why the fast sorts are fast is the main payoff of the
complexity chapter.
\end{ideabox}

\section{3.15 Insertion Sort}

Insertion sort is the first concrete sorting algorithm we'll build.
It's the sort most people \emph{naturally} perform by hand -- think
about how you'd organize a hand of playing cards -- and its
quadratic worst-case cost is the price of its simplicity. It's also
surprisingly good when the input is already nearly sorted, and that
adaptivity is where it shines.

\subsection{The idea}

\begin{defnbox}[Insertion sort]
\textbf{Insertion sort} partitions the array into a \emph{sorted
prefix} and an \emph{unsorted suffix}. On each step, it takes the
next element from the suffix and slides it leftward through the
prefix until it's in the right spot, extending the sorted region by
one.
\end{defnbox}

\begin{ideabox}[Card-sorting analogy]
Imagine picking up cards one at a time. The cards in your hand are
always sorted; each new card from the deck gets walked leftward
until it slots into place. Insertion sort is exactly that, in code.
\end{ideabox}

\subsection{Insertion sort in C++}

\begin{examplebox}[Swap-based insertion sort]
\begin{lstlisting}[basicstyle=\ttfamily\small,frame=single,language=C++]
void insertionSort(std::vector<int>& nums) {
    for (std::size_t i = 1; i < nums.size(); ++i) {
        std::size_t j = i;
        while (j > 0 && nums[j] < nums[j - 1]) {
            std::swap(nums[j], nums[j - 1]);
            --j;
        }
    }
}
\end{lstlisting}
The outer loop marks the boundary between sorted and unsorted. The
inner loop walks \texttt{nums[i]} leftward one swap at a time until
its left neighbor is no larger.
\end{examplebox}

\begin{examplebox}[Shift-based insertion sort (one fewer write per step)]
\begin{lstlisting}[basicstyle=\ttfamily\small,frame=single,language=C++]
void insertionSort(std::vector<int>& nums) {
    for (std::size_t i = 1; i < nums.size(); ++i) {
        int key = nums[i];
        std::size_t j = i;
        while (j > 0 && nums[j - 1] > key) {
            nums[j] = nums[j - 1];   // shift right
            --j;
        }
        nums[j] = key;               // drop key into place
    }
}
\end{lstlisting}
Instead of swapping pairs (three writes each), this version shifts
and then drops the key once. Same algorithm, measurably faster
constant factor.
\end{examplebox}

\begin{notebox}[Two versions, same algorithm]
Swap-based is the textbook version -- it reads more naturally. The
shift-based version is the one you'd write after profiling. Either
is $O(n^2)$ in the worst case; the shift-based version has a smaller
constant.
\end{notebox}

\subsection{Worst-case analysis: $O(n^2)$}

\begin{ideabox}[Counting comparisons]
The outer loop runs $n - 1$ times. On iteration $i$, the inner loop
does at most $i$ comparisons -- it might walk the value all the way
to index $0$. Summing:
\[
  1 + 2 + 3 + \cdots + (n - 1) = \frac{n(n - 1)}{2} \in \Theta(n^2).
\]
This maximum is achieved when the input is in \emph{reverse sorted
order} -- every new element has to travel all the way to the front.
\end{ideabox}

\subsection{Best case: $O(n)$ on already-sorted input}

\begin{ideabox}[Why insertion sort is \emph{adaptive}]
If the input is already sorted, the inner while condition
\texttt{nums[j]\,<\,nums[j - 1]} is false on the very first check of
every iteration. The inner loop does one comparison and terminates.
Total: $n - 1$ comparisons, zero swaps -- $\Theta(n)$.
\end{ideabox}

\subsection*{Nearly sorted lists}

The adaptivity generalizes. Suppose only a constant number $c$ of
elements are out of place. Then
\begin{itemize}[leftmargin=*]
  \item $n - c$ elements each need 1 comparison to stay put.
  \item $c$ elements each need at most $n$ comparisons.
\end{itemize}
Total work: $(n - c) \cdot 1 + c \cdot n = n - c + cn \in O(n)$ (since
$c$ is constant). Insertion sort runs in \emph{linear time} on
nearly-sorted input.

\begin{examplebox}[Which lists are ``nearly sorted''?]
\begin{center}
\footnotesize
\begin{tabular}{@{}lp{7.5cm}@{}}
\hline
$(1, 2, 3, 4, 5, 6)$                                     & Sorted \\
$(1, 2, 3, 4, 5, 10, 6)$                                 & Nearly sorted (one misplaced) \\
$(4, 5, 17, 25, 89, 14)$                                 & Nearly sorted \\
$(6, 5, 4, 3, 2, 1)$                                     & Worst case -- fully reversed \\
$(9, 3, 7, 1, 8, 2)$                                     & Random -- expect $\Theta(n^2)$ \\
\hline
\end{tabular}
\end{center}
\end{examplebox}

\subsection{Why this matters in practice}

\begin{ideabox}[Insertion sort's real-world niche]
Given its $O(n^2)$ worst case, insertion sort sounds like a
pedagogical stop on the way to quicksort. In practice it's still
used, in two specific places:
\begin{itemize}[leftmargin=*]
  \item \textbf{Small inputs.} For $n$ below roughly 10--32, the
        tiny constant factor beats the overhead of fancier sorts.
        Introsort (\texttt{std::sort}) actually switches to
        insertion sort at the leaves of its recursion for exactly
        this reason.
  \item \textbf{Streaming input.} If elements arrive one at a time
        and you want the collection kept sorted, insertion sort is
        what you're doing whether you named it or not.
\end{itemize}
\end{ideabox}

\subsection{Properties}

\begin{center}
\footnotesize
\renewcommand{\arraystretch}{1.25}
\begin{tabular}{@{}ll@{}}
\hline
\textbf{Property}   & \textbf{Insertion sort} \\
\hline
Best case           & $\Theta(n)$ -- already-sorted input \\
Worst case          & $\Theta(n^2)$ -- reverse-sorted input \\
Average case        & $\Theta(n^2)$ for random input \\
Space               & $O(1)$ extra -- in-place \\
Stable?             & Yes -- equal elements retain original order \\
Adaptive?           & Yes -- fast on nearly-sorted input \\
\hline
\end{tabular}
\end{center}

\begin{warnbox}[Stability needs the strict inequality]
The inner test is \texttt{nums[j] < nums[j - 1]}, \emph{not}
\texttt{<=}. Using \texttt{<=} would swap equal elements, which
destroys stability. When you're sorting by a secondary key on top
of a previous sort, stability matters -- use strict \texttt{<}.
\end{warnbox}

\begin{notebox}[Loop invariants: the formal correctness-proof technique]
The intuitive argument ``insertion sort works because each pass
extends the sorted prefix by one'' is a \emph{loop invariant}
proof, just stated informally. CLRS Ch~2.1 makes this rigorous --
the same 3-part structure works for every loop-based algorithm in
this course (binary search, BST insert, Dijkstra, \ldots).

\textbf{The technique.} For insertion sort's outer loop, the
invariant is:

\begin{quote}
\emph{At the start of each iteration, the subarray
\texttt{nums[0..i-1]} contains the original first $i$ elements,
sorted in non-decreasing order.}
\end{quote}

A complete proof has three parts:
\begin{enumerate}[leftmargin=*,noitemsep]
  \item \textbf{Initialization.} The invariant holds before the
        first iteration. Here: \texttt{i = 1}, so the prefix is
        \texttt{nums[0..0]} -- one element, trivially sorted.
  \item \textbf{Maintenance.} If the invariant holds at the start
        of one iteration, it holds at the start of the next. Here:
        the inner while loop slides \texttt{nums[i]} leftward past
        anything larger; when it stops, the prefix
        \texttt{nums[0..i]} is sorted. Increment \texttt{i};
        invariant restored for next iteration.
  \item \textbf{Termination.} When the loop exits, the invariant
        plus the exit condition imply the desired result. Here:
        \texttt{i == nums.size()}, so the entire array is the
        sorted prefix.
\end{enumerate}

That's a complete correctness argument. The same scaffold scales
to harder algorithms -- binary search's invariant is ``the target,
if present, lies in \texttt{[lo, hi]}''; BST search's is ``the
target, if present, lies in the subtree rooted at the current
node.'' Once you can state the invariant, the proof writes itself.
\end{notebox}

\begin{notebox}[Where this is going]
The next sorting sections introduce quicksort (fast average,
worst-case $n^2$ if unlucky with pivots), merge sort (guaranteed
$n \log n$, $O(n)$ space), and heap sort (guaranteed $n \log n$,
in-place). Each earns its place in the toolbox by trading off among
the five properties in the table above.
\end{notebox}

\section{3.16 Selection Sort}

Selection sort is insertion sort's mirror image. Both carve the
array into a sorted prefix and an unsorted suffix; both extend the
prefix by one element per outer-loop pass. The difference is
\emph{which element} gets placed each pass. Insertion sort takes
whatever is next and slides it into position. Selection sort picks
the \emph{smallest element of the remaining unsorted region} and
parks it at the end of the sorted region.

\begin{defnbox}[Selection sort]
\textbf{Selection sort} partitions the array into a sorted prefix
and an unsorted suffix. On each pass, it scans the suffix to
\emph{select} the smallest (or largest, for descending) remaining
element and \emph{swaps it} into the first position of the suffix,
extending the sorted prefix.
\end{defnbox}

\subsection{Selection sort in C++}

\begin{examplebox}[Selection sort]
\begin{lstlisting}[basicstyle=\ttfamily\small,frame=single,language=C++]
void selectionSort(std::vector<int>& nums) {
    for (std::size_t i = 0; i + 1 < nums.size(); ++i) {
        std::size_t indexSmallest = i;
        for (std::size_t j = i + 1; j < nums.size(); ++j) {
            if (nums[j] < nums[indexSmallest]) {
                indexSmallest = j;
            }
        }
        std::swap(nums[i], nums[indexSmallest]);
    }
}
\end{lstlisting}
Outer loop advances the boundary. Inner loop scans the suffix and
records the index of the smallest. One swap per outer iteration
places it.
\end{examplebox}

\begin{warnbox}[Avoid \texttt{i < size - 1} on unsigned types]
\texttt{std::vector::size()} returns \texttt{size\_t} (unsigned). If
the vector is empty, \texttt{size - 1} underflows to a huge number
and the loop runs forever. Write \texttt{i + 1 < nums.size()}
instead -- it's the same condition for non-empty input and safe for
empty input.
\end{warnbox}

\subsection{Why it always does $\Theta(n^2)$ comparisons}

\begin{ideabox}[Counting comparisons]
On pass $i$ (zero-based), the inner loop compares every suffix
element against the current smallest -- that's $n - i - 1$
comparisons. Summing:
\[
  (n - 1) + (n - 2) + \cdots + 1 = \frac{n(n - 1)}{2} \in \Theta(n^2).
\]
The key detail is that the number of comparisons \emph{does not
depend on the data}. Even if the array is already sorted, the inner
loop still scans the entire suffix to \emph{prove} it has found the
minimum. Every case -- best, average, worst -- is $\Theta(n^2)$.
\end{ideabox}

\begin{warnbox}[Selection sort is not adaptive]
Insertion sort runs in $O(n)$ on nearly-sorted input; selection sort
does not. It does the same work no matter what. If your input is
already mostly sorted, reach for insertion sort instead.
\end{warnbox}

\subsection{Where selection sort \emph{is} better than insertion sort}

Selection sort has one redeeming feature: it does at most
\emph{$n - 1$ swaps} total. Insertion sort can do $\Theta(n^2)$
swaps.

\begin{ideabox}[When swap cost dominates]
If a ``swap'' is expensive -- say each element is a large struct
that's costly to move, or is stored in flash memory where writes
wear out the hardware -- selection sort wins. It finds the right
element via cheap comparisons, then does one swap per pass. You
trade comparison work for write work, and sometimes the writes are
what you're actually trying to minimize.
\end{ideabox}

\subsection*{Selection sort vs.\ insertion sort at a glance}

\begin{center}
\footnotesize
\renewcommand{\arraystretch}{1.3}
\begin{tabular}{@{}lll@{}}
\hline
\textbf{Property}     & \textbf{Insertion sort} & \textbf{Selection sort} \\
\hline
Best case             & $\Theta(n)$             & $\Theta(n^2)$ \\
Worst case            & $\Theta(n^2)$           & $\Theta(n^2)$ \\
Space                 & $O(1)$                  & $O(1)$ \\
Stable?               & Yes                     & No (swap can reorder equals) \\
Adaptive?             & Yes                     & No \\
Swaps                 & Up to $\Theta(n^2)$     & Exactly $n - 1$ \\
\hline
\end{tabular}
\end{center}

\begin{warnbox}[Selection sort is \emph{not} stable -- a subtle why]
Consider sorting $[5_a, 3, 5_b, 1]$ (the subscripts distinguish two
``5''s with equal key but different identity). Pass 0 finds the
smallest (1) at index 3 and swaps it with index 0, giving
$[1, 3, 5_b, 5_a]$. The order of the two 5's has flipped even though
they're equal by key. Any secondary ordering you had \emph{is gone}.
That's why when stability matters -- e.g.\ sorting rows by column A
after already sorting by column B -- use insertion sort or merge
sort, not selection sort.
\end{warnbox}

\subsection{When to reach for each}

\begin{ideabox}[A pocket rubric]
\begin{itemize}[leftmargin=*]
  \item \textbf{Insertion sort.} Small $n$, streaming data, nearly
        sorted input, or when stability matters.
  \item \textbf{Selection sort.} Small $n$ where \emph{writes} are
        much more expensive than comparisons.
  \item \textbf{Otherwise.} \texttt{std::sort} (introsort).
\end{itemize}
For real work, the third bullet covers 99\% of cases. The first two
are about understanding the trade-offs so you know \emph{why} the
third bullet is the default.
\end{ideabox}

\begin{notebox}[The pedagogical point]
Insertion and selection sort together illustrate that ``$O(n^2)$''
isn't one algorithm -- it's a class, with members that differ in
\emph{which} operation they do $n^2$ of (comparisons vs swaps),
whether they adapt to input shape, and whether they preserve
equal-element order. Analyzing any new sort means asking the same
checklist.
\end{notebox}

\section{3.17 Shell Sort}

Shell sort is the first sort we meet that breaks out of the
quadratic class -- not by doing something dramatically new, but by
applying insertion sort cleverly in \emph{multiple passes over
interleaved sublists}. It's the bridge between the simple sorts of
3.15--3.16 and the $O(n \log n)$ sorts of the next chapter. Named
for Donald Shell, 1959.

\subsection{The core insight}

Insertion sort is slow on unsorted input because each misplaced
element can only move \emph{one slot at a time}. A tiny value at the
end of a big array has to be swapped past every larger neighbor,
one pair at a time, to reach the front -- that's the $n^2$.

\begin{ideabox}[Move elements in big jumps first]
What if elements could \emph{leap} most of the distance in a few
moves, and only the last little bit was done by ordinary insertion
sort (which is $O(n)$ on nearly-sorted input)? That's Shell sort:
first run insertion sort on many interleaved subsequences with a
big ``gap,'' which lets elements jump far; then shrink the gap
repeatedly; finish with gap = 1 (ordinary insertion sort) on a list
that's already almost sorted.
\end{ideabox}

\subsection{Interleaved lists and the ``gap''}

\begin{defnbox}[Gap]
A \textbf{gap} $K$ is a positive integer. With gap $K$, Shell sort
treats the array as $K$ \emph{interleaved} sublists: indices $0, K,
2K, \ldots$; then $1, K+1, 2K+1, \ldots$; and so on through index
$K-1$. Each sublist is sorted by an insertion-sort variant that
uses $K$ as the step size instead of $1$.
\end{defnbox}

\begin{examplebox}[Gap of 3 on a 9-element list]
Array: \quad $[56, 42, 93, 12, 77, 82, 75, 91, 36]$

With gap 3, the three interleaved sublists are:
\begin{itemize}[leftmargin=*]
  \item indices $\{0, 3, 6\}$: $(56, 12, 75)$
  \item indices $\{1, 4, 7\}$: $(42, 77, 91)$
  \item indices $\{2, 5, 8\}$: $(93, 82, 36)$
\end{itemize}
Each is sorted independently. After sorting the three interleaved
lists, the array looks like:
$[12, 42, 36, 56, 77, 82, 75, 91, 93]$.

The array isn't sorted yet, but it's \emph{close}: the three
smallest elements all ended up in the first three slots. That's the
setup Shell sort creates so that the final gap-1 pass (plain
insertion sort) runs in nearly linear time.
\end{examplebox}

\subsection{Shell sort in C++}

\begin{examplebox}[Shell sort implementation]
\begin{lstlisting}[basicstyle=\ttfamily\small,frame=single,language=C++]
// Insertion sort with a custom step size
static void insertionSortGap(std::vector<int>& nums,
                             std::size_t start,
                             std::size_t gap) {
    for (std::size_t i = start + gap; i < nums.size(); i += gap) {
        std::size_t j = i;
        while (j >= gap && nums[j] < nums[j - gap]) {
            std::swap(nums[j], nums[j - gap]);
            j -= gap;
        }
    }
}

void shellSort(std::vector<int>& nums) {
    for (std::size_t gap = nums.size() / 2; gap > 0; gap /= 2) {
        for (std::size_t start = 0; start < gap; ++start) {
            insertionSortGap(nums, start, gap);
        }
    }
}
\end{lstlisting}
The outer loop shrinks the gap (halving here). For each gap, the
inner loop runs an insertion-sort variant on each of the $K$
interleaved sublists. The final pass uses gap = 1, which is plain
insertion sort.
\end{examplebox}

\begin{warnbox}[You \emph{must} end with a gap of 1]
A Shell sort that never runs with gap = 1 doesn't actually sort.
It just produces an array whose interleaved sublists are sorted.
The final gap = 1 pass is what finishes the job.
\end{warnbox}

\subsection{Choice of gap sequence}

\begin{defnbox}[Gap sequence]
A \textbf{gap sequence} is the descending list of gap values Shell
sort uses, ending with 1. Any sequence works \emph{for correctness},
but the sequence dramatically affects the \emph{runtime}.
\end{defnbox}

\begin{center}
\footnotesize
\renewcommand{\arraystretch}{1.25}
\begin{tabular}{@{}p{4.5cm}p{4cm}p{4.5cm}@{}}
\hline
\textbf{Gap sequence}        & \textbf{Example ($n=100$)}       & \textbf{Worst-case runtime} \\
\hline
Halving: $n/2, n/4, \ldots, 1$ (original Shell) & 50, 25, 12, 6, 3, 1 & $O(n^2)$ \\
Hibbard: $2^k - 1$           & 63, 31, 15, 7, 3, 1              & $O(n^{3/2})$ \\
Sedgewick / Knuth            & 121, 40, 13, 4, 1                & $O(n^{4/3})$ \\
Ciura (empirical, best known small-$n$) & 57, 23, 10, 4, 1        & very fast in practice \\
\hline
\end{tabular}
\end{center}

\begin{ideabox}[The strange fact about Shell sort]
Shell sort's exact asymptotic worst case is still open for some gap
sequences. For Hibbard ($2^k - 1$) it's $O(n^{3/2})$. For
Sedgewick's sequence it's $O(n^{4/3})$. Ciura's sequence isn't known
in closed form but beats the others in benchmarks. Shell sort is
one of the few algorithms in everyday use whose worst-case complexity
depends on an unproven conjecture about a numeric sequence.
\end{ideabox}

\subsection{Why this beats plain insertion sort}

\begin{ideabox}[Why the gap matters]
Plain insertion sort moves misplaced elements \emph{one slot per
swap}. If an element is $d$ slots from its destination, it needs
$d$ swaps. Summed over all elements, that's the $n^2$.

With gap $K$, a swap moves an element $K$ slots. A few gap-$K$
passes can reduce the \emph{total distance every element has to
travel} dramatically before the final gap-1 cleanup. The cleanup is
still insertion sort, but on a list where elements are at most a
small constant away from their destination -- and insertion sort on
nearly-sorted data is $O(n)$.
\end{ideabox}

\subsection{Properties}

\begin{center}
\footnotesize
\renewcommand{\arraystretch}{1.3}
\begin{tabular}{@{}ll@{}}
\hline
\textbf{Property}            & \textbf{Shell sort} \\
\hline
Best case                    & $O(n \log n)$ (already sorted, depends on sequence) \\
Worst case                   & $O(n^2)$ (na\"ive halving) down to $O(n^{4/3})$ (Sedgewick) \\
Average case                 & Depends on gap sequence; between $O(n^{5/4})$ and $O(n^{3/2})$ in practice \\
Space                        & $O(1)$ extra -- in-place \\
Stable?                      & No \\
Adaptive?                    & Somewhat -- fewer passes on presorted input \\
\hline
\end{tabular}
\end{center}

\begin{warnbox}[Shell sort is \emph{not} stable]
When gap $> 1$, swapping non-adjacent elements can reorder equal
keys unpredictably. Use merge sort or insertion sort when stability
is needed.
\end{warnbox}

\subsection{When to use it}

\begin{ideabox}[Shell sort's niche]
Shell sort is rarely the right choice in modern C++ code --
\texttt{std::sort} beats it. But it has a few real uses:
\begin{itemize}[leftmargin=*]
  \item \textbf{Embedded / constrained environments.} In-place,
        simple to implement, no recursion stack, no extra memory.
        Useful when you can't afford $\log n$ stack depth.
  \item \textbf{Educational.} It's the easiest example of ``clever
        insertion sort'' and bridges $n^2$ and $n \log n$.
  \item \textbf{Moderate $n$ in code without \texttt{<algorithm>}.}
        If you can't lean on the STL, Shell sort is a few lines of
        code that outperform plain insertion sort by a lot.
\end{itemize}
\end{ideabox}

\begin{notebox}[Where this is going]
Shell sort showed that \emph{structure in the swap pattern} can
break the $n^2$ barrier. The next step is algorithms that break it
on purpose, with provable $O(n \log n)$ guarantees: quicksort, merge
sort, and heap sort. Shell sort is the last sort where the question
``what gap sequence do I use?'' really matters -- after this, the
algorithms have more definite shapes.
\end{notebox}

\section{3.18 Quicksort}

Quicksort is the first algorithm we meet that achieves
$O(n \log n)$ \emph{on average} -- and it's the most widely used
comparison sort in practice. Its strategy is a textbook example of
\textbf{divide-and-conquer}: do a bit of work to partition the
problem into two smaller subproblems, recurse on each, and stitch
the result together. Invented by Tony Hoare in 1959 while he was a
graduate student trying to sort Russian words for machine
translation.

\subsection{The idea: partition and recurse}

\begin{defnbox}[Quicksort]
\textbf{Quicksort} sorts a range by:
\begin{enumerate}[leftmargin=*]
  \item Choosing a \textbf{pivot} value from the range.
  \item \textbf{Partitioning} the range so every value $\le$ pivot is
        on the left, every value $\ge$ pivot is on the right.
  \item Recursively quicksorting the left side and the right side.
\end{enumerate}
The range is sorted when both recursive calls return.
\end{defnbox}

\begin{ideabox}[Why partitioning works]
Once partitioning is done, \emph{every element on the left is $\le$
every element on the right}. So sorting the two sides independently
sorts the whole range -- no merge step needed. The work to combine
the results is zero. That's what makes quicksort an in-place sort.
\end{ideabox}

\subsection{The Hoare-style partition}

The textbook's partition is a variant of Hoare's original. Two index
pointers walk from the ends toward each other, swapping out-of-place
values when they meet misplaced elements.

\begin{examplebox}[Hoare partition in C++]
\begin{lstlisting}[basicstyle=\ttfamily\small,frame=single,language=C++]
static std::size_t partition(std::vector<int>& nums,
                             std::size_t low,
                             std::size_t high) {
    std::size_t mid = low + (high - low) / 2;
    int pivot = nums[mid];

    while (true) {
        while (nums[low]  < pivot) ++low;
        while (pivot < nums[high]) --high;

        if (low >= high) return high;

        std::swap(nums[low], nums[high]);
        ++low;
        --high;
    }
}

void quicksort(std::vector<int>& nums,
               std::size_t low,
               std::size_t high) {
    if (low >= high) return;                          // base case
    std::size_t mid = partition(nums, low, high);
    quicksort(nums, low, mid);                         // left half
    quicksort(nums, mid + 1, high);                    // right half
}
\end{lstlisting}
Pivot picked at the midpoint; two pointers walk inward; swap when
both find something on the wrong side; stop when they meet. After
the partition returns, each side is recursively sorted.
\end{examplebox}

\begin{warnbox}[Watch the pivot-index caveat]
Hoare's partition returns an index $j$ such that everything in
$[low, j]$ is $\le$ pivot and everything in $[j+1, high]$ is
$\ge$ pivot. The \emph{pivot value itself might end up on either
side}, and there's no guarantee the pivot sits at index $j$. That's
why the recursive calls use \texttt{[low, mid]} and
\texttt{[mid + 1, high]} -- never excluding the middle index, as
you'd do in Lomuto's partition. Getting this wrong is a classic
bug.
\end{warnbox}

\begin{notebox}[Lomuto vs.\ Hoare partition]
Two partition schemes are standard. \textbf{Hoare} (shown above)
uses two converging pointers; it's faster in practice but
\emph{doesn't} place the pivot at its final index. \textbf{Lomuto}
uses a single running ``boundary'' index and pivots around the
\emph{last} element; it's easier to get right, slower by a constant,
and places the pivot at its final sorted index. Intro CS courses
often teach Lomuto; introsort (used in \texttt{std::sort}) uses a
Hoare variant.
\end{notebox}

\subsection{Runtime: average vs.\ worst}

\begin{ideabox}[Average case: $O(n \log n)$]
If the pivot splits the range into two roughly equal halves, the
recursion tree has depth $\log_2 n$ and at each level quicksort does
$\Theta(n)$ work partitioning. Total: $\Theta(n \log n)$.
\[
  \text{levels} \times \text{work per level}
  \;=\; \log n \;\times\; n \;=\; n \log n.
\]
This is what happens on random or ``normal'' data.
\end{ideabox}

\begin{warnbox}[Worst case: $O(n^2)$]
If the pivot is always the smallest (or largest) element, one side
gets zero elements and the other gets $n - 1$. The recursion depth
becomes $n$, and total work is $\Theta(n^2)$:
\[
  n + (n - 1) + (n - 2) + \cdots + 1 = \frac{n(n+1)}{2}.
\]
\textbf{This is especially bad for already-sorted or reverse-sorted
input if the pivot is ``first'' or ``last.''} Na\"ive quicksort on
sorted data is actually \emph{slower than insertion sort}.

\emph{Picking the middle element as pivot} (as the textbook does)
partially mitigates this -- already-sorted inputs partition cleanly.
Adversarial inputs can still force $O(n^2)$, which is why real
implementations use more defensive strategies.
\end{warnbox}

\subsection{Pivot selection strategies}

\begin{center}
\footnotesize
\renewcommand{\arraystretch}{1.3}
\begin{tabular}{@{}p{3.2cm}p{9.2cm}@{}}
\hline
\textbf{Strategy} & \textbf{Behavior} \\
\hline
First or last element   & Disastrously $O(n^2)$ on sorted inputs. Don't. \\
Middle element          & Good for sorted inputs; still $O(n^2)$ on adversarial data. \\
Random element          & Expected $O(n \log n)$ no matter the input pattern.
                          Adversary can't predict the pivot. \\
Median-of-three (first, middle, last) & Cheap approximation of the
                          true median; very robust in practice. The
                          choice most real implementations make. \\
Introsort               & Start with quicksort, but if recursion depth
                          exceeds $2 \log n$, switch to heap sort.
                          Guaranteed $O(n \log n)$ worst case.
                          \texttt{std::sort} does this. \\
\hline
\end{tabular}
\end{center}

\begin{ideabox}[Why \texttt{std::sort} is basically always good enough]
\texttt{std::sort} uses introsort:
\begin{itemize}[leftmargin=*]
  \item Quicksort with median-of-three pivot selection,
  \item Switches to heap sort if recursion is getting too deep
        (guarding the worst case),
  \item Switches to insertion sort on small subranges (where
        insertion sort's small constant wins).
\end{itemize}
You get average quicksort speed, worst-case heap-sort guarantee,
and small-range insertion-sort wins, all in one call.
\end{ideabox}

\subsection{Properties}

\begin{center}
\footnotesize
\renewcommand{\arraystretch}{1.3}
\begin{tabular}{@{}ll@{}}
\hline
\textbf{Property}   & \textbf{Quicksort} \\
\hline
Best case           & $\Theta(n \log n)$ \\
Average case        & $\Theta(n \log n)$ \\
Worst case          & $\Theta(n^2)$ -- mitigated by pivot strategy \\
Space               & $O(\log n)$ recursion stack (in-place for the data) \\
Stable?             & No -- partition swaps non-adjacent elements \\
Adaptive?           & Not especially -- depends on pivot choice \\
\hline
\end{tabular}
\end{center}

\begin{warnbox}[Not stable, again]
If stability matters -- you're sorting by a secondary key on top of
a previous sort -- use \texttt{std::stable\_sort} (merge sort) or
insertion sort, not \texttt{std::sort}.
\end{warnbox}

\subsection{Why is it called \emph{quick} if it can be $O(n^2)$?}

\begin{ideabox}[The constant factor is extremely small]
On random data, quicksort does noticeably fewer operations per
element than merge sort or heap sort, and its inner loop is very
cache-friendly: it streams through memory sequentially and only
touches two array slots at a time. ``Quick'' refers to real
wall-clock performance, not asymptotic analysis. Modern engineered
quicksort (introsort) beats both merge sort and heap sort on typical
inputs, and avoids the $O(n^2)$ pathology by design.
\end{ideabox}

\begin{notebox}[Where this is going]
Quicksort is average-case optimal and the default sort in every
major language. The next algorithm, merge sort, gives up the
in-place property in exchange for a guaranteed $O(n \log n)$ worst
case and stability -- useful when you cannot tolerate the rare
quadratic blowup or when the data is in a linked list (where
partitioning doesn't work but merging does).
\end{notebox}

\section{3.19 Merge Sort}

Quicksort gets its $O(n \log n)$ \emph{on average} and pays a rare
$O(n^2)$ on bad pivots. Merge sort buys a \emph{guaranteed}
$O(n \log n)$ worst case plus stability, at the cost of $O(n)$ extra
memory. It's also divide-and-conquer, but with the work distributed
differently: quicksort does the work \emph{before} recursing
(partition then recurse), merge sort does the work \emph{after}
(recurse then merge). Invented by John von Neumann in 1945.

\subsection{The idea: recurse then merge}

\begin{defnbox}[Merge sort]
\textbf{Merge sort} sorts a range by:
\begin{enumerate}[leftmargin=*]
  \item If the range has $\le 1$ element, it's already sorted.
        Return.
  \item \textbf{Divide} the range into two halves at the midpoint.
  \item Recursively merge-sort each half.
  \item \textbf{Merge} the two sorted halves into a single sorted
        range.
\end{enumerate}
\end{defnbox}

\begin{ideabox}[The two halves of divide-and-conquer]
\textbf{Quicksort} does a smart partition (the hard work) and then
the recursive calls trivially concatenate. \textbf{Merge sort}
divides trivially (just split at the midpoint) and then does the
hard work merging two already-sorted halves. The total work is the
same, $\Theta(n \log n)$, but it lives in different places.
\end{ideabox}

\subsection{The merge step}

The merge is where everything happens. Two already-sorted
subsequences can be combined into one sorted sequence in linear
time by walking both with two pointers and always taking the
smaller of the two front elements.

\begin{examplebox}[Merging two sorted runs in C++]
\begin{lstlisting}[basicstyle=\ttfamily\small,frame=single,language=C++]
static void merge(std::vector<int>& nums,
                  std::size_t i, std::size_t j, std::size_t k) {
    // nums[i..j]   is the left  (sorted) half
    // nums[j+1..k] is the right (sorted) half
    std::vector<int> merged;
    merged.reserve(k - i + 1);

    std::size_t left = i, right = j + 1;
    while (left <= j && right <= k) {
        if (nums[left] <= nums[right]) {
            merged.push_back(nums[left++]);
        } else {
            merged.push_back(nums[right++]);
        }
    }
    while (left  <= j) merged.push_back(nums[left++]);
    while (right <= k) merged.push_back(nums[right++]);

    for (std::size_t p = 0; p < merged.size(); ++p) {
        nums[i + p] = merged[p];
    }
}
\end{lstlisting}
Scan both halves in lockstep. Whichever front element is smaller,
take it. When one half is exhausted, flush the other. Finally, copy
the merged buffer back to \texttt{nums[i..k]}.
\end{examplebox}

\begin{warnbox}[Stability lives in the \texttt{<=}]
The test is \texttt{nums[left] <= nums[right]} (``if tied, take from
the \emph{left}''). This preserves the original order of equal
elements -- that's why merge sort is \textbf{stable}. If you change
it to \texttt{<}, equals from the right get placed first and
stability is lost.
\end{warnbox}

\subsection{Merge sort in C++}

\begin{examplebox}[The recursive driver]
\begin{lstlisting}[basicstyle=\ttfamily\small,frame=single,language=C++]
void mergeSort(std::vector<int>& nums,
               std::size_t i, std::size_t k) {
    if (i >= k) return;                             // base case
    std::size_t j = i + (k - i) / 2;                // midpoint
    mergeSort(nums, i, j);                           // sort left
    mergeSort(nums, j + 1, k);                       // sort right
    merge(nums, i, j, k);                            // combine
}
\end{lstlisting}
The recursion tree has depth $\log_2 n$; each level does $\Theta(n)$
work in its merge calls. Total: $\Theta(n \log n)$.
\end{examplebox}

\subsection{Runtime: why always $O(n \log n)$}

\begin{ideabox}[The recursion tree argument]
Each level of the recursion:
\begin{itemize}[leftmargin=*]
  \item Splits all ranges at their midpoints -- free, just an
        index calculation.
  \item Merges pairs of sorted subranges back together -- each
        element of the original array is touched exactly once by a
        merge at that level, so total work per level is $\Theta(n)$.
\end{itemize}
The tree has depth $\lceil \log_2 n \rceil$. Total work:
\[
    \underbrace{\log_2 n}_{\text{depth}} \;\times\; \underbrace{n}_{\text{work per level}}
    \;=\; \Theta(n \log n).
\]
\textbf{Unlike quicksort, this bound is independent of input order.}
Best case, average case, and worst case are all $\Theta(n \log n)$.
\end{ideabox}

\begin{notebox}[The Master Theorem: a formal tool for divide-and-conquer]
The recursion-tree argument above is one way to solve a recurrence.
For the broad class $T(n) = a \cdot T(n/b) + f(n)$ -- ``solve $a$
subproblems of size $n/b$, plus do $f(n)$ work to combine them'' --
there's a closed-form result called the \textbf{Master Theorem}
(CLRS Ch~4 / OCW r03). Three cases, depending on how $f(n)$
compares to $n^{\log_b a}$ (the work that would happen if you only
counted the leaves of the recursion tree):

\begin{center}
\small
\begin{tabular}{@{}p{2.6cm}p{4.0cm}p{5.0cm}@{}}
\textbf{Case} & \textbf{When} & \textbf{Then $T(n) =$} \\
\hline
1. Leaves win  & $f(n) = O(n^{\log_b a - \varepsilon})$ for some $\varepsilon > 0$ & $\Theta(n^{\log_b a})$ \\
2. Balanced    & $f(n) = \Theta(n^{\log_b a})$                                     & $\Theta(n^{\log_b a} \log n)$ \\
3. Root wins   & $f(n) = \Omega(n^{\log_b a + \varepsilon})$ + regularity         & $\Theta(f(n))$ \\
\end{tabular}
\normalsize
\end{center}

\textbf{Merge sort:} $T(n) = 2T(n/2) + \Theta(n)$. Here $a=2, b=2$,
so $\log_b a = 1$, and $f(n) = n = \Theta(n^1)$ -- case 2 applies.
$T(n) = \Theta(n \log n)$. The hand-derivation above is a worked
instance of case 2. The Master Theorem is just the general formula
that lets you skip the recursion-tree drawing for any recurrence in
this form.

\textbf{Binary search:} $T(n) = T(n/2) + O(1) = 1 \cdot T(n/2) +
\Theta(n^0)$, so $\log_b a = 0$, $f(n) = \Theta(n^0)$ -- case 2
again. $T(n) = \Theta(\log n)$.

\textbf{Quicksort (best case):} $T(n) = 2T(n/2) + \Theta(n)$, same
as merge sort, so $\Theta(n \log n)$. (Worst case requires a
different recurrence and falls outside the Master Theorem.)

A handful of useful recurrences and their solutions:

\begin{center}
\small
\begin{tabular}{@{}ll@{}}
$T(n) = T(n - 1) + O(1)$        & $\Theta(n)$ \\
$T(n) = T(n - 1) + O(n)$        & $\Theta(n^2)$ \\
$T(n) = T(n/2) + O(1)$          & $\Theta(\log n)$ \\
$T(n) = 2T(n/2) + O(1)$         & $\Theta(n)$ \\
$T(n) = 2T(n/2) + O(n)$         & $\Theta(n \log n)$ \\
$T(n) = 2T(n/2) + O(n \log n)$  & $\Theta(n \log^2 n)$ \\
$T(n) = 2T(n - 1) + O(1)$       & $\Theta(2^n)$ \\
\end{tabular}
\normalsize
\end{center}

The full proof of the Master Theorem -- including which $f(n)$ shapes
fall into which case -- is in CLRS Ch~4. For ch\_3 the takeaway is
that the recursion-tree arguments throughout \textsection 3.18 and
\textsection 3.19 are instances of a general result, not one-off
clever tricks.
\end{notebox}

\subsection{Memory cost}

\begin{warnbox}[Merge sort is not in-place]
The merge step writes to a temporary buffer of the same size as the
range it's merging. At the top of the recursion the buffer is size
$n$. That's $O(n)$ extra memory -- dominant over the $O(\log n)$
recursion stack. Quicksort sorts \emph{in place} with only $O(\log n)$
extra memory for the stack; merge sort does not.

For huge datasets that don't fit in RAM, or for embedded systems
where memory is scarce, this can be the deciding factor.
\end{warnbox}

\begin{notebox}[One-time buffer allocation]
A small engineering win: allocate the full-size scratch buffer
\emph{once}, outside the recursion, and pass it to each merge call.
This avoids $\log n$ distinct allocations. The implementation above
reallocates inside \texttt{merge} for clarity; production versions
pass a buffer.
\end{notebox}

\subsection{Properties}

\begin{center}
\footnotesize
\renewcommand{\arraystretch}{1.3}
\begin{tabular}{@{}ll@{}}
\hline
\textbf{Property} & \textbf{Merge sort} \\
\hline
Best case         & $\Theta(n \log n)$ \\
Average case      & $\Theta(n \log n)$ \\
Worst case        & $\Theta(n \log n)$ -- guaranteed \\
Space             & $O(n)$ extra -- temporary merge buffer \\
Stable?           & Yes \\
Adaptive?         & Not natively (plain merge sort). Variants
                    like Timsort are adaptive. \\
\hline
\end{tabular}
\end{center}

\subsection{Merge sort vs.\ quicksort}

\begin{center}
\footnotesize
\renewcommand{\arraystretch}{1.3}
\begin{tabular}{@{}lll@{}}
\hline
\textbf{Criterion}    & \textbf{Quicksort}          & \textbf{Merge sort} \\
\hline
Best / avg case       & $\Theta(n \log n)$          & $\Theta(n \log n)$ \\
Worst case            & $\Theta(n^2)$                & $\Theta(n \log n)$ \\
Extra memory          & $O(\log n)$ (stack)          & $O(n)$ (buffer) \\
Stable?               & No                           & Yes \\
Cache behavior        & Excellent (sequential)       & Good (streaming) \\
In-place?             & Yes                          & No \\
Recursion depth       & $O(\log n)$ avg, $O(n)$ worst & $O(\log n)$ always \\
\hline
\end{tabular}
\end{center}

\begin{ideabox}[When to reach for each]
\begin{itemize}[leftmargin=*]
  \item \textbf{Quicksort} (\texttt{std::sort}) when you have plenty
        of RAM, you don't need stability, and you want typical
        speed.
  \item \textbf{Merge sort} (\texttt{std::stable\_sort}) when you
        need stability or a guaranteed worst case, and you can
        afford the extra memory.
  \item \textbf{Merge sort on linked lists.} Quicksort's partition
        assumes random access. Merge sort's \emph{sequential} access
        pattern makes it the natural choice for linked lists --
        \texttt{std::list::sort} uses it.
  \item \textbf{External sorting (data too large for RAM).} Merge
        sort is the basis for sorting data on disk: sort chunks that
        fit in memory, merge them in passes. Quicksort doesn't
        work that way.
\end{itemize}
\end{ideabox}

\subsection*{Timsort: merge sort's real-world descendant}

\begin{notebox}[Timsort]
Python's built-in \texttt{sorted} and Java's
\texttt{Arrays.sort} for objects use \textbf{Timsort} -- a hybrid
merge/insertion sort designed for real-world data. Timsort detects
already-sorted runs, sorts small ranges with insertion sort, and
merges the runs intelligently. On many real datasets (mostly sorted
with pockets of noise) it beats plain quicksort. It's the most
commonly-run sorting algorithm in the world today.
\end{notebox}

\begin{ideabox}[Where this is going]
Merge sort completes the picture of general-purpose $O(n \log n)$
sorts: in-place but randomized (quicksort), predictable but
memory-hungry (merge sort). Heap sort -- coming later with priority
queues -- is the third option: guaranteed $O(n \log n)$ and
in-place, at the cost of being slower in practice and not stable.
Each point in the trade-space has a winner.
\end{ideabox}

\section{3.20 Radix Sort}

Every sort we've seen so far is \emph{comparison-based}: it uses
\texttt{<} between pairs of elements as its only decision procedure.
The lower bound for that class is $\Omega(n \log n)$. Radix sort is
our first non-comparison sort -- it never compares two elements to
each other -- and it breaks the $n \log n$ bound, sorting in
$O(n \cdot d)$ where $d$ is the number of digits. For fixed-width
integer keys, that's effectively $O(n)$.

\subsection{The trick: sort by one digit at a time}

\begin{defnbox}[Radix sort]
\textbf{Radix sort} sorts an array of integers by processing
\emph{one digit at a time}, from the least-significant digit (LSD)
to the most-significant. For each digit, it distributes values into
10 \textbf{buckets} (one per digit value $0, 1, \ldots, 9$), then
rebuilds the array by concatenating the buckets in order. After
processing all digits, the array is sorted.
\end{defnbox}

\begin{ideabox}[Why this works: stability of the digit pass]
After the first pass, the array is sorted by the 1's digit. After
the second pass -- \emph{sorted by the 10's digit in a stable way} --
ties on the 10's digit keep the order established by the first pass,
so the array is now sorted by ``10's digit, then 1's digit'' --
which is numerical order up to two digits. Repeating for higher
digits extends the sorted-ness to the full number. Radix sort is
correct \emph{only because} the digit-by-digit pass is stable.
\end{ideabox}

\subsection{The LSD algorithm}

\begin{examplebox}[Radix sort in C++, non-negative integers]
\begin{lstlisting}[basicstyle=\ttfamily\small,frame=single,language=C++]
void radixSort(std::vector<int>& nums) {
    if (nums.empty()) return;

    int maxVal = *std::max_element(nums.begin(), nums.end());
    std::vector<std::vector<int>> buckets(10);

    for (int pow10 = 1; maxVal / pow10 > 0; pow10 *= 10) {
        for (int v : nums) {
            int digit = (v / pow10) % 10;
            buckets[digit].push_back(v);
        }
        std::size_t idx = 0;
        for (auto& bucket : buckets) {
            for (int v : bucket) nums[idx++] = v;
            bucket.clear();
        }
    }
}
\end{lstlisting}
Outer loop advances through the digits (ones, tens, hundreds, \ldots).
Inner first loop distributes values to buckets; inner second loop
gathers them back in order. The pass is stable because we iterate
\texttt{nums} in order and \texttt{push\_back} preserves arrival
order within a bucket.
\end{examplebox}

\begin{warnbox}[Stability of the gather step is essential]
If you reorder within a bucket -- for example, sorting or
randomizing each bucket's contents -- radix sort breaks. The whole
algorithm depends on each digit pass being a \emph{stable} partial
sort.
\end{warnbox}

\subsection*{Tracing a small example}

\begin{center}
\footnotesize
\renewcommand{\arraystretch}{1.3}
\begin{tabular}{@{}l|l@{}}
\hline
\textbf{Step} & \textbf{Array} \\
\hline
Input           & 170, 45, 75, 90, 802, 24, 2, 66 \\
After 1's pass  & 170, 90, 802, 2, 24, 45, 75, 66 \\
After 10's pass & 802, 2, 24, 45, 66, 170, 75, 90 \\
After 100's pass & 2, 24, 45, 66, 75, 90, 170, 802 \\
\hline
\end{tabular}
\end{center}

Three passes, because the max value has three digits. After each
pass, the prefix of digit-sorted-ness grows by one.

\subsection{Runtime}

\begin{ideabox}[$O(n \cdot d)$]
Let $n$ be the number of elements and $d$ be the number of digits
of the maximum value. Each pass does $O(n)$ work (distribute +
gather). There are $d$ passes. Total: $\Theta(n \cdot d)$.

For 32-bit integers, $d \le 10$. For 64-bit integers, $d \le 20$.
Those are \emph{constants} -- independent of $n$. So for fixed-width
integers, radix sort is effectively $\Theta(n)$.
\end{ideabox}

\begin{warnbox}[``Beats $n \log n$'' comes with fine print]
Yes, radix sort is linear in $n$ when keys are fixed-size integers.
But:
\begin{itemize}[leftmargin=*]
  \item The constant factor hides ``$\times d$,'' which can be
        larger than $\log_2 n$ for moderate $n$. Radix sort usually
        beats quicksort only for large $n$ with narrow keys.
  \item It doesn't generalize: radix sort only works on things you
        can decompose into digits (integers, fixed-length strings,
        fixed-precision floats with bit manipulation).
  \item It's not in-place in the na\"ive version: 10 buckets of
        potentially $n$ elements each, plus copy-back, is $O(n)$
        extra memory.
\end{itemize}
It's a specialist tool, not a \texttt{std::sort} replacement.
\end{warnbox}

\subsection{Handling negatives and other bases}

\begin{examplebox}[Signed integers: post-process into two buckets]
The simple LSD algorithm above sorts by \emph{absolute value} when
negatives are present (since \texttt{-12} has 1's digit 2). The
textbook's fix: after the main radix sort, split into a
\texttt{negatives} bucket and a \texttt{nonNegatives} bucket,
\emph{reverse} the negatives (so $-42, -12, -78$ becomes
$-78, -42, -12$), and concatenate
\texttt{[reversed negatives] ++ [nonNegatives]}.
\end{examplebox}

\begin{notebox}[Radix = base]
The number of buckets you use is the \emph{radix} or base. Ten is
natural for decimal integers, but you can use any base.
\textbf{Base 256} (one byte at a time) is the usual choice for
radix-sorting binary integer keys on modern computers -- one memory
read per digit, exactly 4 passes for a 32-bit int. Base 2 (one bit
at a time) takes more passes but uses only 2 buckets. Picking the
base is a time / memory trade-off.
\end{notebox}

\subsection{Properties}

\begin{center}
\footnotesize
\renewcommand{\arraystretch}{1.3}
\begin{tabular}{@{}ll@{}}
\hline
\textbf{Property}  & \textbf{Radix sort (LSD)} \\
\hline
Time (best / avg / worst) & $\Theta(n \cdot d)$ -- all cases the same \\
Space              & $\Theta(n + B)$ where $B$ is the base \\
Stable?            & Yes -- and correctness depends on it \\
Adaptive?          & No -- doesn't notice already-sorted input \\
Works on?          & Integer keys (or anything decomposable into
                    fixed-width digits) \\
\hline
\end{tabular}
\end{center}

\subsection{When to reach for radix sort}

\begin{ideabox}[Niche but powerful]
\begin{itemize}[leftmargin=*]
  \item \textbf{Huge volumes of integers or fixed-length strings.}
        Log files with numeric IDs, network packet sorting by IP
        address, database sorts on integer keys -- radix sort can be
        $3$--$10\times$ faster than quicksort.
  \item \textbf{Hardware where comparison is expensive.} GPUs and
        FPGAs often use radix sort because distribution-style
        algorithms parallelize cleanly.
  \item \textbf{Fixed-width keys with a small base.} Sort 8-bit
        values with a single 256-bucket pass -- basically free.
\end{itemize}
For variable-length strings, general objects, or anything that can
only be compared (not decomposed), stick to \texttt{std::sort}.
\end{ideabox}

\subsection{Counting sort: the standalone non-comparison sort}

Radix sort uses counting sort once per digit. But counting sort is
also a complete sorting algorithm in its own right -- the right
choice when keys are integers in a known small range and you want
$O(n + k)$ time, where $k$ is the size of the range.

\begin{defnbox}[Counting sort]
\textbf{Counting sort} sorts integer keys in the range $[0, k)$ by:
(1) counting how many times each key value occurs, (2) computing
cumulative sums to determine each value's final output position,
and (3) walking the input in reverse to place each element at its
position (decrementing the count as we go). $\Theta(n + k)$ time,
$\Theta(n + k)$ space, stable.
\end{defnbox}

\begin{examplebox}[Counting sort in C++]
\begin{lstlisting}[basicstyle=\ttfamily\small,frame=single,language=C++]
void countingSort(std::vector<int>& nums, int k) {
    // Assumes every nums[i] satisfies 0 <= nums[i] < k.
    std::vector<int> count(k, 0);
    for (int v : nums)              ++count[v];          // count occurrences
    for (int i = 1; i < k; ++i)     count[i] += count[i-1]; // cumulative
    std::vector<int> out(nums.size());
    for (int i = nums.size() - 1; i >= 0; --i) {         // walk BACKWARD
        out[--count[nums[i]]] = nums[i];
    }
    nums = std::move(out);
}
\end{lstlisting}
The cumulative-sum array tells you exactly where each value's
``block'' ends in the output. Walking the input \emph{backward} and
decrementing as you go preserves the original order of equal keys
-- that's why counting sort is stable.
\end{examplebox}

\begin{warnbox}[Why backward, why decrement?]
If you walked the input forward and \emph{incremented}, equal keys
would land in their final block in \emph{reverse} input order --
breaking stability. The walk-backward + decrement convention is
what makes counting sort a stable building block, which is exactly
the property radix sort needs.
\end{warnbox}

\begin{ideabox}[When counting sort wins]
\begin{itemize}[leftmargin=*,noitemsep]
  \item \textbf{Grades 0--100.} $k = 101$ buckets, $n$ students.
        Sort by score in $\Theta(n + 101) = \Theta(n)$.
  \item \textbf{ASCII characters.} $k = 128$. Sort a string by
        character in $\Theta(n)$.
  \item \textbf{Pixel intensities, age in years, days of the
        year.} Anywhere keys live in a known small range.
\end{itemize}
Counting sort breaks down when $k$ is huge -- $k = 10^9$ means a
billion-entry array even if $n = 10$. That's where radix sort
re-enters: chop the key into base-$n$ digits and run counting sort
on each digit, getting $\Theta(n)$ for fixed-width integer keys.
\end{ideabox}

\begin{ideabox}[Where this is going]
Radix sort closes out the sorting tour. You now have the
comparison-based trio -- insertion, quicksort, merge sort -- plus
the specialized non-comparison contender, radix sort. The Big-O
machinery you used throughout this chapter (formalized in
\textsection 3.10) is the cost-analysis foundation for every later
chapter. Chapter 4 starts the per-structure walk -- linked lists,
the alternative to contiguous storage -- and every performance
comparison there is phrased in the vocabulary you built here.
\end{ideabox}

\section{3.21 Overview of Fast Sorting Algorithms}

The last several sections introduced a parade of sorting algorithms.
This section steps back and classifies them. Two questions organize
the view: \textbf{how fast?} (average-case complexity) and
\textbf{what kind?} (does the algorithm compare elements, or does it
exploit more structure in the keys?).

\begin{defnbox}[Fast sorting algorithm]
A \textbf{fast sorting algorithm} is one whose average-case runtime
is $O(n \log n)$ or better. Sorts that average $\Theta(n^2)$ -- or
anything slower than $n \log n$ -- are considered ``slow''
regardless of their constant factor.
\end{defnbox}

\subsection{The speed classification}

\begin{center}
\footnotesize
\renewcommand{\arraystretch}{1.3}
\begin{tabular}{@{}lll@{}}
\hline
\textbf{Algorithm} & \textbf{Average case} & \textbf{Fast?} \\
\hline
Selection sort   & $\Theta(n^2)$          & No \\
Insertion sort   & $\Theta(n^2)$          & No \\
Shell sort       & $\Theta(n^{1.5})$ (typical sequences) & No \\
Quicksort        & $\Theta(n \log n)$     & Yes \\
Merge sort       & $\Theta(n \log n)$     & Yes \\
Heap sort        & $\Theta(n \log n)$     & Yes \\
Radix sort       & $\Theta(n \cdot d) = O(n)$ for fixed-width & Yes \\
\hline
\end{tabular}
\end{center}

\begin{ideabox}[The $n \log n$ frontier]
Three comparison sorts meet the $n \log n$ bar on average:
\textbf{quicksort, merge sort, heap sort}. They're not
interchangeable -- they differ on worst case, memory, stability, and
cache behavior -- but they're the short list for ``fast comparison
sort.'' Shell sort sits just below, in the $n^{1.5}$ band;
insertion and selection sort sit well below, in the $n^2$ band.
\end{ideabox}

\subsection{The $\Omega(n \log n)$ lower bound for comparison sorts}

\begin{ideabox}[Why no comparison sort can do better than $n \log n$]
A sorting algorithm that uses only pairwise comparisons is deciding
which of $n!$ possible input permutations it's looking at. Each
comparison has two outcomes, so after $k$ comparisons the algorithm
has distinguished among at most $2^k$ permutations. For the
algorithm to correctly identify any of $n!$ inputs, we need
\[
    2^k \ge n!,
    \quad\text{so}\quad
    k \ge \log_2(n!) \in \Theta(n \log n).
\]
This is a \emph{lower bound on the problem}, not on any particular
algorithm. Quicksort, merge sort, and heap sort all meet it (up to
constants).
\end{ideabox}

\subsection{Comparison vs.\ non-comparison sorts}

\begin{defnbox}[Comparison sort]
A \textbf{comparison sort} uses only the operation ``compare two
elements'' to order them. Every sort except radix sort we've seen is
comparison-based.
\end{defnbox}

\begin{center}
\footnotesize
\renewcommand{\arraystretch}{1.3}
\begin{tabular}{@{}ll@{}}
\hline
\textbf{Algorithm} & \textbf{Comparison sort?} \\
\hline
Selection, insertion, shell sort      & Yes \\
Quicksort, merge sort, heap sort       & Yes \\
Radix sort (and counting, bucket sort) & No \\
\hline
\end{tabular}
\end{center}

\begin{ideabox}[Why radix sort can beat the lower bound]
The $\Omega(n \log n)$ lower bound applies only to comparison sorts.
Radix sort doesn't compare elements; it looks at digits directly,
which is additional structural information. The lower bound doesn't
apply, and radix sort achieves $\Theta(n)$ on fixed-width integer
keys.

The trade-off is universality. A comparison sort works on anything
you can define \texttt{<} for -- strings, structs, floats, custom
types. Radix sort needs keys you can decompose into digits. This is
why \texttt{std::sort} is comparison-based even though radix sort
could be faster for certain inputs: the library has to work on
\emph{any} \texttt{operator<}-able type.
\end{ideabox}

\begin{notebox}[Bucket sort: the third non-comparison sort]
The third member of the non-comparison family is \textbf{bucket
sort}: assume input values are uniformly distributed in $[0, 1)$,
divide the range into $n$ equal-width buckets, drop each input
into its bucket, sort each bucket with insertion sort, and
concatenate. $\Theta(n)$ expected time on uniform input,
$\Theta(n^2)$ worst case. Useful for floating-point keys or other
values you can predict the distribution of. Full coverage in
\textsection 13.3 of the extras chapter.
\end{notebox}

\subsection{Best / average / worst, side by side}

Average case is the headline number, but the other two cases matter
when you care about guarantees (worst) or optimized input (best).

\begin{center}
\footnotesize
\renewcommand{\arraystretch}{1.3}
\begin{tabular}{@{}llll@{}}
\hline
\textbf{Algorithm} & \textbf{Best} & \textbf{Average} & \textbf{Worst} \\
\hline
Quicksort   & $\Theta(n \log n)$  & $\Theta(n \log n)$ & $\Theta(n^2)$ \\
Merge sort  & $\Theta(n \log n)$  & $\Theta(n \log n)$ & $\Theta(n \log n)$ \\
Heap sort   & $\Theta(n)$\,* / $\Theta(n \log n)$ & $\Theta(n \log n)$ & $\Theta(n \log n)$ \\
Radix sort  & $\Theta(n)$         & $\Theta(n)$         & $\Theta(n)$ \\
\hline
\end{tabular}
\end{center}

\begin{notebox}[The $\Theta(n)$ heap-sort best case]
The textbook lists heap sort's best case as $\Theta(n)$. That number
refers to the cost of \emph{building} the heap, which is $\Theta(n)$
by a clever analysis. The subsequent extraction phase is
$\Theta(n \log n)$, so total heap sort work is still
$\Theta(n \log n)$ -- the ``$\Theta(n)$ best case'' applies only to
the heapify stage, not the whole sort. Heap sort the algorithm is
$\Theta(n \log n)$ in all cases.
\end{notebox}

\subsection{Picking a fast sort: a mini decision tree}

\begin{ideabox}[The 80/20 rule]
\begin{enumerate}[leftmargin=*]
  \item \textbf{Just use \texttt{std::sort}.} Introsort. Handles
        almost every case well.
  \item \textbf{Need stability?} \texttt{std::stable\_sort} (merge
        sort). Slightly slower; preserves order of equal keys.
  \item \textbf{Can't afford worst-case blowup?} Merge sort or heap
        sort. Both give $O(n \log n)$ worst case; merge sort uses
        $O(n)$ extra memory, heap sort is in-place.
  \item \textbf{Sorting integers with a narrow range, huge $n$?}
        Radix sort -- the one place $O(n)$ is actually achievable.
  \item \textbf{Data mostly sorted?} Timsort (merge sort variant) or
        insertion sort for tiny $n$. Both exploit existing order.
  \item \textbf{Small $n$ ($< 32$ish)?} Insertion sort wins on
        constants. This is why \texttt{std::sort} falls back to
        insertion sort at the leaves.
\end{enumerate}
\end{ideabox}

\subsection{The extended comparison table}

Everything from the chapter, in one place:

\begin{center}
\footnotesize
\renewcommand{\arraystretch}{1.3}
\begin{tabular}{@{}p{2.2cm}p{1.6cm}p{1.6cm}p{1.6cm}p{1cm}p{1cm}p{1.6cm}@{}}
\hline
\textbf{Algorithm} & \textbf{Best} & \textbf{Average} & \textbf{Worst} & \textbf{Space} & \textbf{Stable?} & \textbf{Adaptive?} \\
\hline
Selection   & $n^2$       & $n^2$       & $n^2$       & $1$   & No  & No  \\
Insertion   & $n$         & $n^2$       & $n^2$       & $1$   & Yes & Yes \\
Shell       & $n \log n$  & $n^{1.5}$   & $n^{1.5\text{--}2}$ & $1$ & No  & Some \\
Quicksort   & $n \log n$  & $n \log n$  & $n^2$       & $\log n$ & No & No  \\
Merge       & $n \log n$  & $n \log n$  & $n \log n$  & $n$   & Yes & No  \\
Heap        & $n \log n$  & $n \log n$  & $n \log n$  & $1$   & No  & No  \\
Radix       & $n$         & $n$         & $n$         & $n$   & Yes & No  \\
\hline
\end{tabular}
\end{center}

\begin{warnbox}[The ``O'' notation hides constants]
Equal Big-O doesn't mean equal speed. On random data, quicksort's
inner loop is two or three times faster than merge sort's --
because quicksort's partition is extremely cache-friendly and
writes only when it has to. ``$n \log n$'' is the right first-order
answer, but the second-order constants decide which $n \log n$
algorithm actually wins on your hardware with your data.
\end{warnbox}

\begin{ideabox}[Where this is going]
This closes the sorting tour. The tools you built here -- Big-O, the
growth hierarchy, and a feel for comparison vs.\ non-comparison sorts --
are load-bearing for every remaining chapter. Linked-list search in
ch.~4 is linear time (and you'll know why without thinking). BST search
in ch.~6 is $O(\log n)$ amortized but $O(n)$ worst-case (and you'll know
which shape of tree triggers which). Hash table lookup in ch.~5 is $O(1)$
expected (and you'll know what ``expected'' is buying you).
\end{ideabox}

\section*{What this connects to}

\begin{notebox}[Forward links to later chapters]
\begin{itemize}
    \item \textbf{Chapter 4 (linked lists).} Every performance comparison
    in ch.~4 is ``compared to a vector, \ldots''. The complexity
    vocabulary from 3.9 and 3.10 is what lets those comparisons land.
    Linked-list traversal is $O(n)$ because each step is constant and
    there are $n$ of them -- the exact reasoning you practiced here.
    \item \textbf{Chapter 5 (hash tables).} The big claim -- ``hash
    lookup is $O(1)$ expected'' -- only makes sense if you're fluent
    with expected vs.\ worst case, which is a 3.10 skill. The chaining
    and probing strategies there are extensions of the collision idea
    introduced at the end of 3.1.
    \item \textbf{Chapter 6 (BSTs).} BST insert/search/remove are
    $O(h)$, which is $O(\log n)$ on a balanced tree and $O(n)$ on a
    degenerate one. The ``balanced'' piece is literally about whether
    the tree shape keeps you on the log curve from 3.10's growth
    hierarchy.
    \item \textbf{Sorting choices in later work.} Whenever you see
    \texttt{std::sort} (introspection sort, essentially quicksort +
    heapsort fallback) you're looking at the quicksort analysis from
    3.18. When you see \texttt{std::stable\_sort}, that's merge sort
    from 3.19 with its stability guarantee.
\end{itemize}
\end{notebox}

\textbf{Companion materials.} Compact notes:
\texttt{chapters/ch\_3/notes.tex} (includes the full sort comparison
table). Practice prompts (problem $\to$ pseudocode $\to$ C++):
\texttt{chapters/ch\_3/practice.md}.

\end{document}
